\section{Proof of Theorem~\ref{thm:main_johnson_close_to_1}}
\label{sec:johnson_1}
In this section we prove Theorem~\ref{thm:main_johnson} 
in the case that $c = 1-\eps$ where $\eps>0$ is small 
so as to isolate the ``conditioning on an event'' challenge as explained in the introduction. In the next section, we explain the modifications that are necessary to prove Theorem~\ref{thm:main_johnson} in full generality. 

We begin by formally defining the Johnson graph.

\begin{definition}[Johnson Graph] \label{def:johnson}
For any $\alpha \in (0,1)$ and $n,\ell \in \N$ with $\alpha\ell \in \N$ and $n > \ell$, we define the \emph{$(n,\ell,\alpha)$-Johnson graph $J(n,\ell,\alpha\ell )$} to be the graph whose vertex set is $\binom{[n]}{\ell}$ and where edges are between pairs of vertices $U,V \in \binom{[n]}{\ell}$ if and only if $|U \cap V| = (1-\alpha)\ell$. We will drop the $(n,\ell,\alpha\ell)$ when clear from context,  
and use $[N]$ to denote the set of vertices of $J$. 

We will often refer to $\alpha$ as the noise parameter 
of the graph.
\end{definition}

\paragraph{Notation:}
The $(n,\ell,\alpha\ell)$-Johnson graph contains other Johnson graphs as subgraphs, corresponding to the basic sets from 
the introduction. Indeed, for any $S \subset [n]$ of size 
smaller than $\ell$ we may consider the induced subgraph on 
all vertices $A$ such that $A\supseteq S$. We will often 
refer to such subgraphs as $|S|$-restrictions and denote it by $J|_S$. The 
motivation for this name is that for given a function $F\colon V(J)\to \mathbb{R}$ 
and $S$, we can define the restricted function 
$F|_{S}\colon V(J|_S)\to\mathbb{R}$ by $F|_{S}[A] = F[A\cup S]$. We will use the notation $\delta(F|_S)$ to denote $\E_{u \in J|_S}[F(u)]$.

% \begin{definition}[$r$-restricted subcubes of $J$]\label{def:subcube}
% Given an $(n,\ell,\alpha)$-Johnson graph $J$ and a set $a \subseteq [n]$ with $|a| = r$ such that $0 \leq r \leq \ell-1$, we let $J|_{a}$ denote the vertex-induced subgraph of $J$ induced by vertices that contain the set $a$. We call such a subset an $r$-restricted subcube of $J$. Note that when $a = \emptyset$ and $r = 0$, $J|_a$ is defined as the whole graph $J$.
% \end{definition}


\subsection{The Algorithm for Affine Unique Games over the Johnson Graph}
In this section, we describe our algorithm for Affine Unique-Games on $J(n,\ell,\alpha)$. Throguhout, when 
we say $q = \exp(r)$ we mean that there exist universal positive constants $c_1, c_2$ such that $2^{c_1 r} \leq q \leq 2^{c_2 r}$. Further if we set a parameter $q$ to be $\exp(r)$ we mean that we set $q = 2^{cr}$ for some universal constant $c$. We also have
an error parameter $\nu>0$ which we determine
in the end but should be thought of as small but bounded away from $0$.

\paragraph{High level description of the algorithm}
\begin{enumerate}
    \item First, our algorithm solves the degree $D$ sum of squares relaxation of program~\eqref{eq:ip} to find a pseudo-expectation $\pE_{\mu}$ corresponding to a pseudo-distribution 
    whose value is at least $1-\eps$.
    \item Next, using the manipulations that were described above, we produce a new one corresponding to two independent samples, 
    as in Sections~\ref{sec:ind_samples}.
    \item In the new pseudo-distribution, it holds that the expected objective function for both the samples $X$ and $X'$ is at least $1-\eps$, hence morally the fraction of edges that are simultaneously satisfied by $X$ and $X'$ is $1-2\eps$. We argue that in this case, there exists a subcube $C$ corresponding to a restriction of constant size on 
    which we have a large \emph{shift-partition potential} \footnote{The existence 
    of such a subcube is argued via global hypercontractivity,
    and to ensure this applies for our pseudodistribution 
    we have to make sure that the global hypercontractivity 
    result we used is proved within the sum-of-squares 
    system.}. We formally define this notion soon, but remark for now that it guarantees that a simple rounding procedure as described in the introduction manages to satisfy a constant fraction of the constraints inside the subcube.
    \item We use a subprocedure, that we refer to as SubRound, to find a subcube $C$ with large shift-partition potential and then an $\Omega(1)$-valued assignment for the vertices in the subcube $C$.
    %We defer the formal description of this procedure to 
    % a later point, and for now it will be sufficient for us 
    % to know that this procedure is given a pseudo-expectation 
    % $\pE$ and a subcube $C$ in which the value of $\pE$ 
    % is high, and in return it finds an assignment $f\colon C\to\Sigma$ satisfying many of the edges inside $C$.
    \item We randomize the edges incident on the vertices of $C$ to get a new instance $I'$ (with possibly lower value) 
    and iterate the algorithm again. The goal of this step is 
    that in the next iteration, we will find a subcube $C'$ 
    which only has very small overlap with the vertices we have assigned thus far. Satisfying a large fraction of $C'$ in the same manner as we did for $C$ would then ensure that we have satisfied new edges and made progress. Iterating thus, we manage to satisfy an $\Omega(1)$-fraction of the graph. This kind of iteration was also used in~\cite{BBKSS}, but at each step they only managed to satisfy an $\Omega(1/\ell^r)$-fraction of the constraints in an $r$-subcube.
\end{enumerate}

\subsubsection{Shift-partition Potential~\cite{BBKSS}}
Our analysis uses the definition of the shift-partition potential from~\cite{BBKSS}, but we explain it here for completeness. 
Recall that given two assignments $X$ and 
$X'$ to an affine Unique-Games instance $I$, 
one may define the shift partition 
$F_s$ which consists of vertices $v$ for 
which $X(v) - X'(v) = s$. As explained 
in the introduction we will heavily use the shift partition, but for 
technical reasons we only want to work with vertices $u$ on which many of the constraints adjacent to them are satisfied. For an assignment $X$ and a vertex $v \in V(G)$, let ${\sf val}_v^G(X)$ denote the value of $v$, i.e. the fraction of the edges of $G$ that are incident on $v$ and are also satisfied by $X$. We drop the superscript $G$ when clear from context.

Naturally, 
this means that we want to include $v$ in $F_s$ only if $X(v) - X'(v) = s$ and ${\sf val}^G_v(X), {\sf val}^G_v(X')$ are at least 
somewhat large. To stay within the realm of
\sos though, we replace these indicator constraints with approximating polynomials 
as in Theorem~\ref{thm:step-approx}.


\begin{definition}[Shift-Partition~\cite{BBKSS}]
Given two affine UG assignments $X$ and $X'$ to $G$ define the following functions $\{F_s:V(G) \to \R\}_{s \in \Sigma}$ \begin{equation}
F_{s}^G(u) = \Ind(X_u - X'_u = s) \cdot p_{\beta,\nu}(\val^{G}_u(X))p_{\beta,\nu}(\val_u^{G}(X')),\label{eq:fs-j}
\end{equation}
where $p_{\beta,\nu}(x)$ is the degree-$\tilde{O}(1/\nu)$ polynomial which is a polynomial approximation to $\Ind[x \geq \beta]$ with accuracy parameter $\nu$ as in Theorem~\ref{thm:step-approx}. We drop the superscript $G$ in $F_s^G$ when the graph is clear from context.
\end{definition}

% If the $p$'s were actual indicator functions instead of polynomial-approximations to them, the functions $F_s$ partition the graph covering almost all the vertices. 
% In that case, any edge that goes across parts in the shift-partition must be violated by either $X$ or $X'$. Hence if $X,X'$ were $(1-\eps)$-satisfying assignments the shift-partition is non-expanding, in the sense that at most $2\eps$-fraction of the edges go across the parts. 

% For small-set expanders,~\cite{BBKSS} then conclude that the fractional-size of one of the parts must be at least $\Omega(\delta)$, and in fact the quantity $\sum_s \E_{u \sim G}[F_s(u)]^2$ must be $\geq \Omega(\delta)$.
% On the other hand they prove that if the squared-size is large on average over $(X,X') \sim \mu \times \mu$ then Condition\&Round on $\mu$ succeeds in rounding to an $\Omega_{\eps,\delta}(1)$-assignment. To aid the analysis they define the shift-partition potential as follows:

Equipped with the definition of shift partition, we can now define the shift-partition potential that governs 
our most basic rounding procedures inside 
subcubes.
\begin{definition}[Shift-Partition Potential~\cite{BBKSS}]\label{def:shift-partition}
The shift-partition size given two UG assignments $X$ and $X'$ on a graph $G$ is defined as:
\[\Phi^G_{\beta,\nu}(X,X') = \sum_{s \in \Sigma} \E_{u \in G}[F_s^G(u)]^2.\]
The shift-partition potential for a pseudodistribution $\cD$ over $(X,X')$ is defined as:
\[\Phi^G_{\beta,\nu}(\cD) = \pE_{(X,X') \sim \cD}[\Phi^G_{\beta,\nu}(X,X')].\]
\end{definition}
Intuitively, $\E_{u \in G}[F_s^G(u)]$ measures the 
fractional size of $F_s$ in the shift partition, 
hence the shift partition potential measures the 
collision probability of the shift partition. In
particular, it can be $\Omega(1)$ if and inly if there
is $s$ such that the part $F_s$ in the partition 
has constant density.

Throughout our arguments, we will consider the shift-partition potential with respect 
to the whole Johnson graph, as well as with 
respect to subcubes inside it. When we consider the shift-partition potential with respect to a subcube $J|_a$ we abbreviate it as $\Phi^a_{\beta,\nu}(\cD)$.





\subsubsection{The Main Routine}
We now give a formal description of our algorithm, modulo the procedure SubRound which we present later.
\begin{algorithm-thm}[Unique Games on the Johnson Graph]\label{alg:j} 
The input to the algorithm is an instance $I$ of affine Unique-Games over a Johnson graph $J(n,\ell,\alpha)$ such that 
${\sf val}(I) \geq 1-\eps$. The output of the algorithm is 
an assignment to $V(J)$.
\begin{enumerate}
\item Let $r =\left\lfloor\frac{64\sqrt{\eta}}{\alpha}\right\rfloor$, $\delta := \frac{\eps}{\exp(r)}$, $\gamma = \eps$ and $D = \frac{|\Sigma|^3\ell^{O(r)}}{\eps^3}$.
Fix $\cA_I$ to be the set of unique games axioms/ integer program over the instance $I$ (Program~\eqref{eq:ip}).

\item Set $j = 1, I_0 = I, R_0 = \phi$.
\item While $|R_{j-1}| < \frac{\gamma}{2}|V(J)|$ do the following:
\begin{compactenum}
\item Solve the degree-$D$ \sos SDP relaxation for the integer program $\cA_{I_{j-1}}$ and apply the transformation 
from Sections~\ref{sec:ind_samples} to get a pseudo-expectation corresponding to two independent samples. Let $\mu_j \times \mu_j$ be the corresponding pseudo-distribution.
\item Find $r' \le r$ and an $r'$-restricted subcube $C_j$ with high SubRound-value\footnote{This is the value of the assignment returned by the subroutine SubRound in~\eqref{alg:j-partial}.}, namely such that $\text{SubRound-val}_{\mu_j}(C_j) \ge \delta$. 

\item Run SubRound (Algorithm~\ref{alg:j-partial}) on $C_j$ with pseudodistribution $\mu_{j}$ to get an assignment $f_j$ to $V(C_j)$.

\item Let $S_j$ be a subgraph of $C_j$ induced by the set of vertices that have not been previously assigned by any partial assignment $f_{k}, k < j$ and assign them using $f_j$. Set $R_j = R_{j-1} \cup H_j$. 

\item Choose random affine-constraints on edges incident on $V(R_j)$ and let the new instance by $I_j$. 

\item Increment $j$.
\end{compactenum}
\item Output any assignment to $V(J)$ that agrees with all partial assignments $f_j$ (assigned to $S_j$) considered above.
\end{enumerate}
\end{algorithm-thm}

\subsubsection{The Subroutine SubRound and Condition\&Round}
Before describing the subroutine SubRound, 
we present several notations that are necessary for the analysis. For a vertex $u\in V(G)$ we will introduce new auxiliary Boolean variables $p_u$ and $p'_u$.  Morally speaking, $p_u$ are the indicators
that many of the constraints adjacent to $u$ are satisfied by the assignment $X$, but formally 
we approximate it by polynomials. 

First note that given a valid degree $D$ pseudodistribution $\cD$ one can extract valid and consistent local distributions over any $D$ variables from $X$, for e.g. looking at the marginal of $(X_u,X_v)$ in $\cD$ we get a distribution over $\Sigma \times \Sigma$. Similarly for any indicator variables $R$ in $\cD$ (where $R_u = R_u^2$ under the axioms of $\cD$) one could analogously define local distributions over $D$ variables from $X,X',R,R'$. Unfortunately the indicator that the value of $u$ is large is not a polynomial. Nevertheless using $\cD$ we can extract the following collection of local distributions over $X,X',p,p'$, where $p$ is a new set of variables we introduce and are not present in $\mu$. We denote this collection by $\cL(\cD, p_{\beta,\nu})$.

Formally, given a pseudodistribution $\cD$ of degree $D > \tO(1/\nu)$ over $(X,X')$ assignments to $\cA_I$ so that the corresponding pseudo-expectation has a high value, 
we define the local distributions $\cL(\cD, p_{\beta,\nu})$ over the variables of $p,p', X, X'$ in the following way:
\begin{definition}[Collection of Local Distributions $\cL(\cD,p_{\beta,\nu})$]\label{defn:local-dist}
Let $\cD$ be a pseudodistribution of degree $D\geq \tilde{\Omega}(d/\nu)$ over $(X,X')$ assignments to $\cA_{I}$, and let the polynomial $p_{\beta,\nu}$ be from Theorem~\ref{thm:step-approx} for $\eps=\delta = \nu$ 
and $\alpha = \beta$.
We define joint distributions over collections of $d$ variables from $p,p',X,X'$ as:
\[\pPr_{{\cL}}[p_u = 1] := \pE_{\cD}[p_{\beta,\nu}(\val_u(X))], ~~~ \pPr_{\cL}[p_u = 0] = 1- \pE_{\cD}[p_{\beta,\nu}(\val_u(X))],~~~ \pPr_{\cL}[X_u = \sigma] = \pE[X_{u,\sigma}],\]
and analogously for $X'_u,p'_u$. More generally we can extend the above definition to define the probability of conjunction of events in the variables $X,X',p,p'$. For all subsets $S,T,A,B \subseteq [N]$ with $|S|+|T|+|A|+|B| \leq d$, $b_1,\ldots,b_{|S|},b'_1,\ldots,b'_{|T|} \in \{0,1\}$ and $\sigma_1,\ldots,\sigma_{|A|},\sigma'_1,\ldots,\sigma'_{|B|} \in \Sigma$ define:
\begin{align*}
&\pPr_{{\cL}}\left[\bigcap_{u \in S} p_u = b_u, \bigcap_{v \in T} p'_v = b'_v, \bigcap_{a \in A} X_{a} = \sigma_a,\bigcap_{b \in B} X'_{b} = \sigma'_b\right] \\= &\pE_{\cD}\left[\prod_{u \in S} (1-b_u + (-1)^{1-b_u}p_{\beta,\nu}(\val_u(X))) \prod_{v \in T} (1-b_v + (-1)^{1-b_v}p_{\beta,\nu}(\val_v(X'))) \prod_{a \in A}X_{a,\sigma_a} \prod_{b \in B}X'_{b,\sigma_b}\right].
\end{align*}
\end{definition}
\begin{remark}
Note that since $\cD$ has large enough degree and $\cA_I \vdash_{\tO(1/\nu)} p_{\beta,\nu}(\val_u(X)) \in [0,1]$, it is easy to check that the above collection of local distributions form valid probability distributions over $d$ variables at a time, and are consistent with each other.
\end{remark}

Using the variables $p,p'$ we further define auxiliary random variables $Y_{u,v}$ and $Y'_{u,v}$ for $u \neq v \in [N]$ that will be particularly useful in the description and analysis of our algorithm: 
\begin{equation}
Y_{u,v} = (X_u,X_v,p_u,p_v), ~~~~ Y'_{u,v} = (X'_u,X'_v,p'_u,p'_v).
\end{equation}
Additionally, analogous to the notation $Y_{u,v}$, for all $S = \{i_1,\ldots,i_{|S|}\} \subseteq [N]$, define the random variables $Y_S, Y'_S$ as: 
\[Y_S = (X_{i_1}, \ldots, X_{i_{|S|}}, p_{i_1}, \ldots, p_{i_{|S|}}), ~~~ Y'_S = (X'_{i_1}, \ldots, X'_{i_{|S|}}, p'_{i_1}, \ldots, p'_{i_{|S|}}).\]
From Definition~\ref{defn:local-dist} we know that the event $Y_S = y_S$ naturally corresponds to an \sos polynomial $Q_S$ such that $\pPr_{{\cL}(\mu \times \mu, p_{\beta,\nu})}[Y_S = y_S] = \pE_{\mu}[Q_S]$ and similarly for events on $Y'_S$. Therefore we will use the notation ``condition $\mu$ on $Y_S = y_S$'' to mean reweighting $\mu$ on the corresponding polynomial $Q_S$. We will frequently look at quantities like $\pPr_{\cL}[\Ind[Y_S = y_S] \wedge \Ind[Y_T = y_T]]$ and $I(Y_S;Y_T)$ where $|S \cup T| \leq d$. In such places we emphasize that we're thinking of the local distribution on $Y_{S \cup T}$ and computing probabilities or mutual information on this distribution.

\skipi
\paragraph{The Subroutine Condition\&Round.}
To formally state SubRound, we first need to present
the subroutine Condition\&Round from~\cite{BBKSS} which is similar to 
the basic rounding procedure from the introduction.
% This subroutine is given a shift symmetric pseudo-distribution 
% $\mu$, and produces a good assignment for it provided that the pseduo-distribution is 
% close to a product distribution in some local sense.
\begin{algorithm-thm}[Condition\&Round~\cite{BBKSS}]\label{alg:low-ent-restated} ~\\ 
\textbf{Input:} A degree-$D$ (for $D \geq 2$) 
shift-symmetric pseudodistribution
%\footnote{Any pseudodistribution can be efficiently transformed into a shift-symmetric one without losing value. See Definition~\ref{def:sym} and Lemma~\ref{lem:sym-basic} for details.} 
$\mu$ for a UG  instance $I'=(G=(V,E),\Pi)$ over alphabet $\Sigma$.  \\
\textbf{Output:} Returns an assignment $x \in \Sigma^V$. %satisfying $\Omega(1)$ fraction of the constraints in expectation.
\\\\
Sample a random solution $Z$:
\begin{compactenum}
\item Sample a vertex $u \sim \pi$ and condition on $X_{u} = 0$ to obtain the new marginals $\pE_\mu[\cdot ~|~ X_{u} = 0]$.
\item Sample a solution $Z$ by choosing each collapsed variable's labels independently according to its marginals: $Z_v \sim \pE_\mu[X_v ~|~ X_{u} = 0]$.
\end{compactenum}
\end{algorithm-thm}
\paragraph{Derandomized Condition\&Round.}
For future reference, it will be more convenient for us to analyze
a derandomized version of Condition\&Round, which makes the 
above procedure deterministic using standard methods such as 
the method of conditional expectations. We omit the straightforward details, and for us it suffices that this derandomized 
procedure always achieves value which is at least the 
expected value of what Condition\&Round gives.
\skipi
We can now describe SubRound that obtains a partial assignment for the Johnson graph. 

\paragraph{High level description of SubRound.}
In this routine, we are given a pseudoexpectation and a 
subcube $J|_a$ (on which we hope the value is large), and we wish to 
find a good assignment to the nodes in the given subcube, 
or quit. This method successfully assigns values to $V(J|_a)$ 
in the case that the shift-partition potential function on the subcube is large.

Towards this end, we find an appropriate conditioning on the random variables $Y_{u,v}$ so that after performing it, 
they are nearly pairwise independent (we use mutual information 
to denote this; this is the step that we referred to in 
the introduction as eliminating global correlations). Denote the new distribution obtained as $\mu_1 \times \mu_2$.

Simultaneously we find an event 
(which we state 
as a polynomial to stay within the realm of \sos) $P_a(X,X')$, which roughly says that 
the part $\{v~:~X(v)-X'(v) = s\}$ is dense inside $J|_a$, 
such that this event has significant probability 
and conditioned on it the shift-partition inside $J|_a$
is large, where both these facts are measured with respect to $\mu_1 \times \mu_2$. We give a formal definition of $P_a(X,X')$ in Lemma~\ref{lem:sp-j}.

Finally, we run the procedure Condition\&Round using either the symmetrized versions of $\mu_1$ or $\mu_2$ which 
returns a good assignment to the given subcube provided that the previous steps succeeded.

\skipi
We now move on to the formal description of SubRound.
\begin{algorithm-thm}[SubRound: Rounding for a subgraph]\label{alg:j-partial} ~\\
\textbf{Input:} Takes as input an affine Unique-Games instance $I = (G, \Pi, \Sigma)$ over an $(n,\ell,\alpha\ell)$-Johnson graph $G$ and alphabet $\Sigma$, a pseudo-distribution $\mu$ with $\val_\mu(I) \geq 1- \eta$, a $j$-restricted subcube $J|_a$, with $j \leq r = \left\lfloor\frac{32\sqrt{\eta}}{\alpha}\right\rfloor$. 

\noindent\textbf{Output:} Returns an assignment in $\Sigma^{|V(J|_a)|}$ to the vertices of $J|_a$. 
\begin{enumerate}
\item Fix $t = \frac{|\Sigma|^3\ell^{\Theta(r)}}{\eta^3}, \tau = \Theta(\frac{\eta}{|\Sigma|^2\ell^{2r}\exp(r)}), \beta = \sqrt{\eta}, \nu = \eta\exp(-r)$.
\item Find a polynomial $P_a(X,X')$ from the set of polynomials given by Lemma~\ref{lem:sp-j}, subsets $A, B \subseteq V(J|_a)$ of size at most $t$ and strings $y_A, y_B$ so that conditioning $\mu$ on the events $Y_A = y_A$ and $Y_B = y_B$ gives pseudodistributions $\mu_1, \mu_2$ satisfying:
\begin{enumerate}
\item 
Shift partition potential inside $J|_{a}$ is significant: 
\[
\pE_{\mu_1 \times \mu_2}[\Phi^a_{\beta,\nu}(X,X')P_a(X,X')] \geq \exp(-r) \pE[P_a(X,X')].
\]
\item The probability of $P_a$ is significant:
\[
\pE_{\mu_1 \times \mu_2}[P_a(X,X')] \geq \Omega\left(\frac{\sqrt{\eta}}{|\Sigma|\ell^{j}\exp(r)}\right).
\]
\item The random variables $Y_{u,v}$ have small global correlation inside $J|_a$: 
\[\E_{\substack{u_1,v_1 \sim S: u_1 \neq v_1\\ u_2,v_2 \sim S: u_2 \neq v_2}}[I(Y_{u_1,v_1};Y_{u_2,v_2})] \leq \tau.
\]
\item The random variables $Y_{u,v}'$ have small global correlation inside $J|_a$: \[
\E_{\substack{u_1,v_1 \sim S: u_1 \neq v_1\\ u_2,v_2 \sim S: u_2 \neq v_2}}[I(Y'_{u_1,v_1};Y'_{u_2,v_2})] \leq \tau.
\]
\end{enumerate}
Here $S = V(J|_a)$, and the mutual information is taken with respect to the collection of local distributions $\cL(\mu_1 \times \mu_2,p_{\beta,\nu})$. If no such polynomial or conditioning exists then output the all zeros assignment on $V(J|_a)$ and exit.
\item Symmetrize $\mu_1, \mu_2$ using the procedure in Section~\ref{sec:symm-prelims} to obtain $\mu^{\sym}_1, \mu^{\sym}_2$. Perform derandomized Condition\&Round on $\mu^{\sym}_1$ or $\mu^{\sym}_2$ and output the higher valued assignment obtained for $V(J|_a)$.
\end{enumerate}
\end{algorithm-thm}



\subsection{Analysis of Algorithm~\ref{alg:j}}
\subsubsection{High Level Description of the Analysis}
We prove that Algorithm~\ref{alg:j} returns a solution with value independent of the alphabet size. We begin by explaining the ideas of the analysis~\cite{BBKSS} for small-set expanders, and then 
explain how this analysis is adapted to our case.

\paragraph{Small-set expanders~\cite{BBKSS}:} 
Consider the shift-partition potential, 
and suppose for simplicity of presentation that the $p$'s are replaced with the actual indicator functions. 
Then the functions $F_s$ cover 
almost all the vertices of the graph, 
and any edge that goes across parts in the shift-partition must be violated by either $X$ or $X'$. Hence if $X,X'$ were $(1-\eps)$-satisfying assignments then the shift-partition is non-expanding, in the sense that at most $2\eps$-fraction of the edges go across the parts (those that are violated 
by either $X$ or $X'$). 

In the case of $(2\eps,\delta)$-small-set expanders,~\cite{BBKSS} then conclude that the fractional-size of one of the parts must be at least $\Omega(\delta)$, and in fact the shift-partition size: $\sum_s \E_{u \sim G}[F_s(u)]^2$, must be $\geq \Omega(\delta)$. On the other hand they prove that if the shift-partition potential is large on $\mu \times \mu$ then Condition\&Round on $\mu$ succeeds in rounding to an $\Omega_{\eps,\delta}(1)$-assignment, 
and formally they show:
\begin{enumerate}
\item Given a degree $D$ certificate of $(O(\eps),\delta)$-small-set expansion for a graph $G$ and a degree $\Omega(D)$ pseudodistribution $\mu$ with value $1-\eps$ on any affine UG instance $I = (G,\Pi)$, conclude that $\Phi^G_{\beta,\nu}(\mu \times \mu) \geq \delta$.
\item Show that if $\Phi^G_{\beta,\nu}(\mu \times \mu)$ is large then Condition\&Round outputs an assignment with large value.
\end{enumerate}

\paragraph{Our analysis for Johnson graphs:} 
Let us see how to adapt this analysis to the case of Johnson graphs. Step (1) is far from being true since Johnson graphs are not small-set expanders. But we know that all non-expanding sets in these graphs must have large size when restricted to some subcube, although the set might not be large in the whole graph. We use this property to conclude that the shift-partition potential must be large on a subcube. This part of the analysis turns out to be much more non-trivial than the certifiable-SSE case and is where the bulk of our technical work lies. 

Once we have this we use a suitably modified version of step (2) to conclude that Condition\&Round when applied to the subcube outputs an $\Omega(1)$-assignment to the subcube. We then use an iteration lemma to run this algorithm multiple times to output an $\Omega(1)$-satisfying assignment to the whole graph. 

% In the next few sections, we elaborate on the above steps.



\subsubsection{Lemmas to be Proven Later}
In this section, we state a few lemmas that are necessary for the analysis of 
Algorithm~\ref{alg:j}, whose proofs are deferred to later sections. In Section~\ref{sec:combine_analysis} we 
use these lemmas to analyze the performance 
of Algorithm~\ref{alg:j}.
\skipi

The analysis of Algorithm~\ref{alg:j} proceeds by showing that given a pseudodistribution for $I$ with high value, there exists a subcube with high SubRound value, i.e. showing the success of Step 3(a). To do so, we first prove that the structure theorem for Johnson graphs (Theorem~\ref{thm:structure-johnson}) implies that given a pseudodistribution $\mu$ with large value, there exists a subcube $J|_a$ with large shift-partition potential denoted by $\Phi^a_{\beta,\nu}$. Since the Johnson graph is not a small-set expander though, we have the following more subtle statement: There exists an $r = O(\sqrt{\eps}/\alpha)$-restricted subcube $a$ and an event $P_a$ that has large probability, such that if we condition $\mu \times \mu$ on $P_a$, the induced shift-partition potential is large on the subcube $J|_a$. The event $P_a$ roughly corresponds to the indicator that the size of one of the shift-partition components is large inside the subcube $J|_a$, i.e. $\Ind(\delta(F_s|_a) > \exp(-r))$ and has probability $\geq \ell^{-O(r)}$.  


%The shift-partition potential measures the average (over $X,X' \sim \cD$) size of the partition defined by the functions $F_s$, where $F_s(u) = \Ind(X_u - X'_u = s)\Ind(\val_u(X) > \beta)\Ind(\val_u(X') > \beta)$. Of course since we cannot use indicator functions we use smooth approximations to them represented by  $p_{\beta,\nu}(\val_u(X))$.~\cite{BBKSS} proved that given a high-value pseudodistribution the shift-partition potential must be large on the whole graph. 

%Additionally to aid our rounding analysis later on, we reduce the global correlation of $\mu$ to get pseudodistributions $\mu_1,\mu_2$ such that they satisfy the same as above: $P_a$ has probability at least $\ell^{-O(r)}$ and the induced shift-partition potential on $J|_a$ is large with respect to the distribution $\mu_1 \times \mu_2$ conditioned on $P_a$. As explained earlier, by ``reducing global correlation'' we mean that we make the variables $Y_{u,v}$ to become almost pairwise independent on average over $u,v \sim V(J|_a)$. Formally we prove the following lemma in Section~\ref{sec:large-potential}:


\begin{lemma}\label{lem:sp-j}
There exists a constant $\eps_0 \in (0,1)$, such that for all positive constants $\eps \leq \eps_0$, $\alpha, \tau \leq 1$, $\nu \leq \eps \exp(-r)$, all integers $\ell \geq \Omega(r), n \geq \ell$, where $r = \lfloor \frac{32\sqrt{\eps}}{\alpha}\rfloor$ the following holds. 

Let $I$ be an affine UG instance on $J(n,\ell,\alpha\ell)$ and $\mu$ be a pseudodistribution over assignments for $I$ with $\val_\mu(I) \geq 1-\eps$ and degree at least $D = \widetilde{\Omega}\left(\frac{|\Sigma|\ell^r\exp(r)}{\nu\tau\sqrt{\eps}}\right)$. Then there exists a restriction $a \subseteq [n]$ of size $j \leq r$, a degree $\widetilde{O}(1/\nu)$ polynomial $P_a(X,X')$ from a fixed set of $|\Sigma|$ polynomials, subsets $A, B \subseteq V(J|_a)$ of size at most $\tO(\frac{|\Sigma|\ell^j\exp(r)}{\tau\sqrt{\eps}})$ and strings $y_A, y_B$ such that conditioning $\mu$ on the events $Y_A = y_A$ and $Y_B = y_B$ gives degree $\widetilde{\Omega}(1/\nu)$ pseudodistributions $\mu_1$ and $\mu_2$ such that:
\begin{enumerate}
\item $\cA_I  ~~\vdash_{\widetilde{O}(1/\nu)} ~~ P_a(X,X') \in [0,1]$.
\item $\pE_{\mu_1 \times \mu_2}[\Phi^a_{\sqrt{\eps},\nu}(X,X')P_a(X,X')] \geq \exp(-r) \pE[P_a(X,X')]$.
\item $\pE_{\mu_1 \times \mu_2}[P_a(X,X')] \geq \Omega\left(\frac{\sqrt{\eps}}{|\Sigma|\ell^{j}\exp(r)}\right)$.
\item $\E_{\substack{u_1,v_1 \sim S: u_1 \neq v_1\\ u_2,v_2 \sim S: u_2 \neq v_2}}[I(Y_{u_1,v_1};Y_{u_2,v_2})] \leq \tau.$
\item $\E_{\substack{u_1,v_1 \sim S: u_1 \neq v_1\\ u_2,v_2 \sim S: u_2 \neq v_2}}[I(Y'_{u_1,v_1};Y'_{u_2,v_2})] \leq \tau,$
\end{enumerate}
where $S = V(J|_a)$, $Y_{u,v} = (X_u,X_v,p_u,p_v)$, $Y'_{u,v} = (X'_u,X'_v,p'_u,p'_v)$ and the mutual information is taken with respect to the collection of local distributions $\cL(\mu_1 \times \mu_2,p_{\sqrt{\eps},\nu})$. 
\end{lemma}

% We have a subcube $J|_a$ with high shift-potential when conditioned on $P_a$, and we next show that this implies that $a$ has high SubRound value, i.e. Condition\&Round when applied on $\mu_1$ or $\mu_2$ should succeed in finding a high-valued assignment on $J|_a$.~\cite{BBKSS} showed that if a graph has large 
% shift-partition potential with respect to $\mu \times \mu$, then Condition\&Round on $\mu$ succeeds in obtaining a highly-satisfying assignment. This analysis crucially used the independence of $X, X'$. We need to work with the conditioned distribution $\mu_1 \times \mu_2 | P_a$ though, which might have correlations between $X$ and $X'$. Nevertheless we show that if $E$ is any event that occurs with large probability,
% and the variables $Y_{u,v}$ and $Y'_{u,v}$ are almost pairwise-independent on average, then conditioning on $E$ cannot correlate too many pairs of variables $(Y_a,Y'_a)$. Such a pair was independent if $(X,X') \sim \mu_1 \times \mu_2$, and we show that the distribution of an average such pair after conditioning on $E$ remains close in TV-distance to its original distribution. We have the following lemma:

The next lemma captures the intuition that 
given a collection of random variables $Z_{u,v}$ that are roughly pairwise independent (that are functions of assignments $X$ and $X'$), and an event $E(X,X')$ with significant 
probability, then typically $Z_{u,v}~|~E$ 
is close to $Z_{u,v}$ in statistical distance. Informally, this asserts that 
conditioning on the event $E$ does not 
change the distributions of $(X, X')$ locally.

\begin{lemma}\label{lem:correlation}
For all $\tau,p,\beta,\nu,\delta \in (0,1)$, integers $D,N$, $m \leq N$ the following holds: Let $\mu_1 \times \mu_2$ be a degree $D+\widetilde{\Omega}(1/\nu)$ pseudodistribution over $(X,X')$ satisfying $\cA_I(X,X')$ and $E(X,X')$ be a polynomial such that, $\cA_I \vdash_D E(X,X') \in [0,1]$ and $\pE_{\mu_1 \times \mu_2}[E(X,X')] \geq p$. Suppose for $S \subseteq [N]$ we have that, 
\begin{enumerate}
\item $\E_{\substack{u_1,v_1 \sim S: u_1 \neq v_1\\ u_2,v_2 \sim S: u_2 \neq v_2}}[I(Y_{u_1,v_1};Y_{u_2,v_2})] \leq \tau.$
\item $\E_{\substack{u_1,v_1 \sim S: u_1 \neq v_1\\ u_2,v_2 \sim S: u_2 \neq v_2}}[I(Y'_{u_1,v_1};Y'_{u_2,v_2})] \leq \tau$,
\end{enumerate}
where $Y_{u,v} = (X_u,X_v,p_u,p_v), Y'_{u,v} = (X'_u,X'_v,p'_u,p'_v)$ and the mutual information is with respect to the collection of local distributions $\cL(\mu_1 \times \mu_2,p_{\beta,\nu})$ (Definition~\ref{defn:local-dist}). Then we have that, 
\[\Pr_{u,v \sim S: u \neq v}[ TV((Y_{u,v},Y'_{u,v})|E , (Y_{u,v},Y'_{u,v})) \geq \delta] \leq O\left(\frac{\sqrt{\tau}+ 1/|S|}{p\delta^2}\right),\]
where the distribution $(Y_{u,v},Y'_{u,v})|E$ refers to the joint distribution on these variables defined by the collection of local distributions $\cL(\mu_1 \times \mu_2|E, p_{\beta,\nu})$ and similarly $(Y_{u,v},Y'_{u,v})$ refers to the distribution defined by $\cL(\mu_1 \times \mu_2,p_{\beta,\nu})$.
\end{lemma}


%In the case of certifiable small-set expanders~\cite{BBKSS} proved that if $G$ has large shift-potential then Condition\&Round on $G$ outputs an assignment with large value. The analysis used the independence of $X$ and $X'$ crucially. 

We now restrict our attention to the nice subcube $J|_a$ obtained from Lemma~\ref{lem:sp-j} and we henceforth only care about the relevant pseudodistributions (e.g. $\mu_1 \times \mu_2|P_a$) when marginalized to $J|_a$. For simplicity of notation we still refer to the marginalized pseudodistributions using their original notation.

We show that large shift-potential $\Phi^a(\mu_1 \times \mu_2 | P_a) \geq \delta$ with the additional property that an average $(Y_{u,v}, Y'_{u,v})$-pair when drawn from $\mu_1 \times \mu_2 |P_a$ is close to its distribution when drawn from $\mu_1 \times \mu_2$, implies that Condition\&Round (applied on $J|_a$) succeeds on either $\mu_1$ or $\mu_2$.

\begin{lemma}\label{lem:round-j}
Let $I = (G, \Pi)$ be an affine instance of Unique Games over the alphabet $\Sigma$ and $\mu_1 \times \mu_2$ be a degree $D+\widetilde{O}(1/\nu)$ pseudodistribution over assignments $(X,X')$ to $I$. Let $E(X,X')$ be a polynomial such that $\cA_I(X,X') \vdash_D E(X,X') \in [0,1]$. Suppose we have that: 
\[\Pr_{u,v \sim V(G): u \neq v}[TV((Y_{u,v},Y'_{u,v})|E , (Y_{u,v},Y'_{u,v})) > \delta] \leq \zeta,\]
where the distribution $(Y_{u,v},Y'_{u,v})|E$ refers to the joint distribution on these variables defined by the collection of local distributions $\cL(\mu_1 \times \mu_2|E, p_{\beta,\nu})$ and similarly $(Y_{u,v},Y'_{u,v})$ refers to the distribution defined by $\cL(\mu_1 \times \mu_2,p_{\beta,\nu})$ (Definition~\ref{defn:local-dist}). 

If $\Phi^G_{\beta,\nu}(\mu_1 \times \mu_2 | E(X,X')) \geq \gamma$, then on at least one of the pseudodistributions $\mu^\sym_1$ or $\mu^\sym_2$ Algorithm~\ref{alg:low-ent-restated} runs in time $\poly(|V(G)|)$ and returns an assignment of expected value at least 
\[
(\beta - \nu)^2\left(\gamma - O(\delta+\zeta) - \frac{1}{|V(G)|}\right) - 3\nu(\beta - \nu)
\]
for $I$.
\end{lemma}

Using the lemmas above for $G = J|_a$ one can analyze a single iteration of SubRound and show that it manages to assign a subcube of $I$ and within it satisfy a constant fraction of the edges. This subcube might be of size $o(1)$ of the whole graph though, therefore to complete the analysis of our algorithm we need the following lemma, asserting that we can iterate this procedure (as done in Algorithm~\ref{alg:j}) to satisfy a constant fraction of constraints of $I$. 
% This completes the analysis of SubRound and shows that there exists a subcube with large SubRound value. We are still not done since SubRound only produces an assignment for a $o(1)$-fraction of vertices. We now apply an iteration lemma that converts a partial assignment into an assignment for an $\Omega(1)$ fraction of the vertices of the graph. This lemma uses the fact that SubRound finds an assignment that satisfies a large fraction of the internal edges of a subcube, and the expansion of this subcube is small.

\begin{lemma}\label{lem:subroutine}
Let $c,\gamma,\delta \in (0,1]$ and $r \in \nn$. Let $I$ be an affine UG instance on alphabet $\Sigma$ on $J_{n,\ell,\alpha}$ with $\ell, n$ large enough  and value at least $c$. 
%Let $\delta_{\min} = \min_{\eta \in [c-\gamma,c]}(\delta(\eta))$.
Suppose we have a subroutine $\cA$ which given as input any affine UG instance $I'$ on $J_{n,\ell,\alpha}$ with $\val_\mu(I') \ge c-\gamma$, returns an $\leq r$-restricted subcube $H$ on $J$ and a partial assignment $f$ such that, $\text{val}_{f}(H) \geq \delta$.
% = \delta(c-\val_\mu(I))
Then if $\cA$ runs in time $T(\cA)$, there is a $O(|V(J)|T(\cA)+|V(J)|^3)$-time algorithm which finds a solution for $I$ that satisfies an $\Omega(\delta\gamma(1-\alpha)^r)$-fraction of the  edges of $J$. 
\end{lemma}
\begin{proof}
Deferred to Section~\ref{sec:pf_it}.
\end{proof}

% \begin{lemma}\label{lem:subroutine}
% Let $c,\gamma,\phi \in (0,1]$, let $\delta: [0,1] \rightarrow [0,1]$ be a function and denote $\delta_{\min} = \min_{\eta \in [c-\gamma,c]}(\delta(\eta))$. 
% Let $G$ be a regular graph and $I$ be any unique games instance on $G$ with alphabet size $|\Sigma| = k \geq \Omega(\frac{1}{\delta_{\min}})$ and value at least $c$. 

% Suppose we have a subroutine $\cA$ which given as input $\mu$, a shift-symmetric degree-$D$ pseudodistribution satisfying $\cA_I$ with $\val_\mu(I) \ge c-\gamma$, returns a vertex-induced subgraph $H$ such that: 
% \begin{enumerate}
%     \item $\text{subRound-val}_{\mu}(H) \geq 
%     \delta(\val_\mu(I))$,
%     \item The edge-expansion of $H$ is at most $\phi$. 
% \end{enumerate} 
% Then, there is an algorithm whose run-time is ${\sf time}(\mathcal{A}) \cdot |V(G)|^{O(D)}$ that given an instance $I$ of unique games over $G$ finds 
% an assignment satisfying at least $\Omega(\delta_{\min}^2\gamma(1-\phi))$-fraction of the constraints. 
% \end{lemma}

% This completes the analysis overview of the algorithm. We combine the above lemmas with the correct parameter settings in the next section. The proofs of the lemmas appear in the following sections after that.



\subsubsection{The Analysis of Algorithm~\ref{alg:j}}\label{sec:combine_analysis}

\begin{theorem}\label{thm:main-johnson}
There exists a constant $\eps_0 \in (0,1)$, such that for all positive constants $\eps \leq \eps_0$, $\alpha \in \Q$ with $\alpha \leq 1/2$, all integers $\ell \geq \Omega(r)$ with $\alpha \ell \in \N$ and $r = \Theta(\sqrt{\eps}/\alpha)$ and $n$ large enough, Algorithm~\ref{alg:j} has the following guarantee: If $I$ is an instance of affine Unique Games on the $(n,\ell,\alpha)$-Johnson graph $G$ with alphabet $\Sigma$ and $\val(I) = 1 - \eps$, then in time $|V(J)|^{\poly(\ell^r,|\Sigma|,1/\eps)}$ 
Algorithm~\ref{alg:j} returns an $\frac{\poly(\eps)}{\exp(r)}$-satisfying assignment for $I$.
\end{theorem}

\begin{proof}
Let $r = \lfloor 32\eps/\alpha\rfloor$. Given the $\leq r$-restricted subcube $a \subseteq [n], P_a(X,X')$ and $\mu_1 \times \mu_2$ from Lemma~\ref{lem:sp-j} we have that $\Phi_{\sqrt{\eps},\nu}^a(\mu_1 \times \mu_2| P_a) \geq \exp(-r)$, therefore we apply Lemma~\ref{lem:correlation} with $S = J|_a$, the polynomial  $E(X,X') = P_a(X,X')$ and the parameters $\beta = \sqrt{\eps}, \nu = \exp(-r)\eps, \delta = \exp(-r), p = \Omega\left(\frac{\sqrt{\eps}}{|\Sigma|\ell^{r}\exp(r)}\right)$ and $\tau = \Theta(\frac{\eps}{|\Sigma|^2\ell^{2r}\exp(r)})$. The parameter $\tau$ has been chosen so that $O(\frac{\sqrt{\tau}+1/|S|}{p\delta^2}) = \exp(-r)$ and therefore we can apply the rounding lemma (Lemma~\ref{lem:round-j}) with $\zeta = \exp(-r)$ and $\gamma = \exp(-r)$ and the same settings of $\beta, \nu, \delta$. This shows that there exists a subcube $a \in {[n] \choose j}$ with SubRound value that is at least $\Omega(\eps/\exp(r))$ if $\deg(\mu) \geq \widetilde{\Omega}\left(\frac{|\Sigma|\ell^r\exp(r)}{\nu\tau\sqrt{\eps}}\right)$. Hence it suffices to have degree of $\mu$ equal to $\frac{|\Sigma|^3\ell^{O(r)}}{\eps^3}$.

By using SubRound as a subroutine, we finish the proof of this theorem by applying the iteration Lemma~\ref{lem:subroutine} with $c=1-\eps$, $\gamma = \eps$, $\delta = \Omega(\eps/\exp(r))$ and $r = \lfloor 64\sqrt{\eps}/\alpha \rfloor$ so that $(1-\alpha)^r = O(\sqrt{\eps})$ (Claim~\ref{claim:subcube-expanse}). We can check that SubRound satisfies the hypotheses of the lemma: it finds a $\le r$-restricted subcube and an assignment to it with value $\Omega(\delta)$ in time $|V(J)|^{\frac{|\Sigma|^3\ell^{O(r)}}{\eps^3}}$ (this follows from the degree upper bound on $\mu$).
This gives us that Algorithm~\ref{alg:j} outputs an assignment of value at least $\Omega(\eps^3/\exp(r))$ in time $|V(G)|^{\frac{|\Sigma|^3\ell^{O(r)}}{\eps^3}}$.
\end{proof}



\subsection{Proof of Lemma~\ref{lem:sp-j}: Finding Subcube with Large Shift-partition}\label{sec:large-potential}
%Recall the definition of the shift-partition potential of a pseudodistribution $\cD$ over UG assignments $(X,X')$\footnote{Until now we worked with $\cD = \mu \times \mu$, but from now on $\cD$ might be any joint pseudodistribution over $(X,X')$.}:
%\[\Phi_{\beta,\nu}(\cD) = \pE_{(X,X')\sim \cD}[\ldots].\]

In this section we use the structure theorem for Johnson graphs
to prove the existence of a restriction $a \subseteq [n]$ of constant size (corresponding to the induced subgraph $J|_a$) with large induced shift-partition potential when conditioned on a non-negative polynomial $P_a$, i.e. $\Phi^a_{\beta,\nu}(\mu \times \mu | P_a)$ is large. Recall that to analyze our final Condition\&Round step though we need to show conditioning $\mu \times \mu$ on $P_a$ does not introduce too many correlations between $X, X'$. Therefore we first perform a global correlation reduction procedure on $\mu \times \mu$ to get the product pseudodistribution $\mu_1 \times \mu_2$ while preserving the property that $\Phi^a_{\beta,\nu}(\mu_1 \times \mu_2 | P_a)$ is large.
The exact quantity that aids our rounding analysis is: 
\[I(Y_{u_1,v_1};Y_{u_2,v_2}) =  I(X_{u_1},X_{v_1},p_{u_1},p_{v_1};X_{u_2},X_{v_2},p_{u_2},p_{v_2}),\] 
defined according to the collection of local distributions in Definition~\ref{defn:local-dist}. We upper bound this by using the global correlation reduction procedure of Raghavendra-Tan~\cite{RT12}. 

\subsubsection{Restricting the Shift Partition to a Subgraph}
To prove Lemma~\ref{lem:sp-j} we will need the notion of
\emph{global shift-partition potential restricted to subcube $J|_a$}, which is almost the same as shift-partition potential on $J|_a$ except for one key difference:
\begin{definition}[Global shift-potential restricted to Subgraphs~\cite{BBKSS}]\label{def:res-shift}
For any $\nu,\beta \in (0,1)$ and subgraph $H$ of $G$, define the {\em global shift-partition size restricted to the subgraph $H$} to be the quantity:
\[
\Phi_{\beta,\nu}(X,X')|_H = \sum_{s \in \Sigma} \E_{u \in H} [F_s^{G}(u)]^2,
\]
where $F_s$'s are the functions defining the shift-partition (Definition~\ref{def:shift-partition}). 
Let the global shift-partition potential with respect to a pseudodistribution $\cD$ over pairs of assignments $(X,X')$ to $G$ be:
\[
\Phi_{\beta,\nu}(\cD)|_H = \pE_{(X,X') \sim \cD}[\Phi_{\beta,\nu}(X,X')|_H].
\]
\end{definition}

Note the difference between the global potential and the shift-partition potential on $H$: the global shift-partition potential $\Phi_{\beta,\nu}(X,X')|_{H}$ measures the size of the global partition inside $H$, i.e. $\val_u^G(X)$ is a function of all the edges of the graph that are incident on $u$, not just the edges in the subgraph $H$. The potential 
$\Phi_{\beta,\nu}^{H}(X,X')$, on 
the other hand, measures the value of a vertex $u$ only \emph{inside} the subgraph $H$. 


% We will use the structure theorem for Johnson graphs to get that there exists a subcube $J|_a$ such that one can condition the pseudodistribution on an event $P_a$ such that the global shift-partition potential restricted to $J|_a$, $\Phi(\mu_1 \times \mu_2 | P_a)|_{a}$ is large. The proof of this is analogous to Lemma 4.4\dor{what?} for certifiable small-set expanders. The main difference is in the conclusion of the lemma: instead of getting a lower bound on the shift-partition potential of the whole graph, we get a lower bound on the global shift-partition potential restricted to some subcube (with respect to the conditioned pseudodistribution $\mu \times \mu | P_a$). We then relate the global shift-partition potential to the induced potential on $C$. 
%Properties 4,5 on the mutual information are obtained to aid with the later rounding step on $C$ (Section~\ref{sec:low-ent-j}). 
\subsubsection{Global Hypercontractivity}
We will need the following definition to describe the structure theorem for Johnson graphs:
\begin{definition}[Restrictions of Functions]
For the $(n,\ell,\alpha)$-Johnson graph $J$, given a function $F:V(J) \rightarrow \R$ and a set $a \subseteq [n]$ with $a = r$, such that $0 \leq r \leq \ell-1$, we define the restricted function $F|_a: \binom{[n] \setminus a}{\ell - r} \rightarrow \R$ as, 
$$F|_a(X) = F(a \cup X).$$
Further, let $\delta(F|_a)$ denote the fractional size of the function restricted to the subcube $J|_a$, that is,
$$\delta(F|_a) := \E\limits_{X \sim \binom{[n] \setminus a}{\ell - r}}[F|_a(X)].$$
When $a = \phi$ and $r = 0$, we have that $F|_a(X) = F(X)$ for all $X \in \binom{[n]}{\ell}$ and $\delta(F|_a) = \delta(F) = \E[F]$.
\end{definition}


In~\cite{KMMS} it is shown that pseudorandom sets expand. Formally, we define pseudorandom sets and in general pseudorandom functions as follows:
\begin{definition}[Pseudorandom functions]\label{def:p.r.}
A set $S$ is called $(r,\gamma)$-pseudorandom if for all $\leq r$-restrictions $a$, $\delta(S|_a) \leq \gamma$. Similarly a function $F$ bounded in $[0,1]$ is called $(r,\gamma)$-pseudorandom if for all $\leq r$-restrictions $a$, $\delta(F^2|_a) \leq \gamma$.
\end{definition}

We will need a version of this 
result for general Boolean and pseudorandom functions $F$, and moreover to show that it is proved in the 
\sos proof system of constant degree. We get an \sos proof that shows that if $\delta(F^2|_a) \leq \gamma$ for all $a$ of size $\leq r$ and $F$ is Boolean, then $\ip{F,LF}$ is large, with $q_a(F)$ being the SoS multipliers of the axioms $\{\gamma - \delta(F^2|_a)\}_a$ and $(\Pi_{\geq \lambda_r}F)(X)$ being the multiplier for the axioms $\{F(X)^3 = F(X)\}$. We have the following statement: 

% In~\cite{KMMS} it is shown that if $F$ is a Boolean function which is the 
% indicator of a set $\mathcal{F}\subseteq \binom{[n]}{\ell}$ whose edge 
% expansion is at most $1-\eps$, then $\delta(F|_{a})\geq \Omega(1)$ for 
% some subcube $a$ of constant size. We will need a version of this 
% result for bounded functions, and moreover to show that it is proved in the 
% \sos proof system of constant degree, as follows.

\begin{theorem}[Expansion Theorem for Johnson Graphs]\label{thm:structure-johnson}
For all $\alpha \in (0,1)$, all integers $\ell \geq 1/\alpha$ and $n \geq \ell$, the following holds: Let $J(n,\ell,\alpha \ell)$ be the $\alpha$-noisy Johnson graph. For every constant $\gamma \in (0,1)$ and positive integer $r \leq O(\ell)$, every function $F: V(J) \rightarrow \R$ that is $(r,\gamma)$-pseudorandom has high expansion:
\begin{align}
&\{F(X) \in [0,1]\}_{X \in V(J)}  \vdash_{O(1)}\, \notag\\
&\ip{F,LF} \geq \delta(F)(1 - (1-\alpha)^{r+1})(1- \gamma^{1/3}\exp(r)) - \sum_{j = 0}^r \frac{ c_j\ell^{j}}{\gamma}\E_{a \sim {[n] \choose j}}[q_a(F)({\delta(F^{2}|_a) - \gamma)}]+ B(F) \\
&\{F(X) \in [0,1]\} \,\,\, \vdash_{2} \,\,\, 0 \leq q_a(F) \leq 1 
\end{align}
where for all $j \leq r$, $c_j$'s are positive constants of size at most
$\exp(r)$, for all size $j$ subsets $a \subseteq [n]$, $q_a(F)$ are degree $2$  polynomials and $B(F)=\frac{4}{3}\E_X[(F^{3} - F)\Pi_{\geq \lambda_r}F ]$ for $\lambda_r = (1-\alpha)^r$ and $\Pi_{\geq \lambda}$ denoting the projection operator to the top-eigenspace of $J$ of eigenvalues $\geq \lambda$.
\end{theorem}
\begin{proof}
Deferred to Section~\ref{sec:fourier}.
\end{proof}


One can equivalently view the expansion result in~\cite{KMMS} as asserting that if $F$ is a Boolean function which is the
indicator of a set $\mathcal{F}\subseteq \binom{[n]}{\ell}$ whose edge 
expansion is bounded away from $1$, then $\delta(F|_{a})\geq \Omega(1)$ for 
some $r$-restricted subcube $a$ with constant $r$. In fact, if the expansion of $F$ is at most $1-\eps$ on $J_{n,\ell,\alpha}$, then $\delta(F|_a) \geq \exp(-r)$ for $r = O(\eps/\alpha)$. Similarly, we can use the above \sos statement to conclude that even if $F$ is ``almost-Boolean'' and non-expanding, then $q_a(F)(\delta(F^2|_a) - \gamma) > 0$ for some $O(1)$-restriction $a$. In the context of a pseudodistribution $\mu$ over non-expanding sets $F$, one can conclude that conditioning $\mu$ on $q_a(F)$ results in a new pseudodistribution $\mu'$ where $\delta(F^2|_a) > \gamma$ on average, and therefore $q_a(F)$ roughly corresponds to the ``event'' that $F$ is dense on $J|_a$.

We use the observation that the shift-partition defined with respect to $(1-\eps)$-satisfying assignments $X$ and $X'$ has expansion at most $2\eps$, therefore is non-expanding. Therefore as above, using the structure theorem we conclude that at least one of the sets in the shift-partition is not $(r,\gamma)$-pseudorandom for $r = O(\eps/\alpha)$ and $\gamma = \exp(-r)$, i.e. $\delta(F_s|_a) \geq \gamma$ for some $s \in \Sigma$ and $\leq r$-restriction $a$. Let $\viol(X)$ denote the fraction of edges that an assignment $X$ violates in the instance $I$. We frame this fact in \sos using the following lemma:

\begin{lemma}\label{lem:structure-fs-j}
Under the conditions of Lemma~\ref{lem:sp-j}, for all $\eta > 0$ and $r \leq \ell$ we get:
\begin{align}
\A_I  
&\vdash_{\tilde{O}(1/\nu)} \\ 
&\sum_{j = 0}^r c_j \exp(r) \ell^j \E_{a \in {[n] \choose j}}[\sum_{s}q_{a}(F_s)(\delta(F_s|_a) - \exp(-r))] \notag\\
&\geq \frac{1 - \lambda_{r+1}}{2} 
- (\viol(X)+\viol(X'))\left(1+\frac{3(1-\lambda_{r+1})+ 12 +8\eta}{6(1-\beta-\nu)}\right) - \frac{2}{3\eta} - \nu\left(5 - \lambda_{r+1} + \frac{8\eta}{3}\right)\label{eq:sp-j-1},
\end{align}
where $\lambda_{r+1} = (1-\alpha)^{r+1}$, $q_a$ are the polynomials in Theorem~\ref{thm:structure-johnson}, $\A_I$ is the set of axioms defined for $I$ by program (\ref{eq:ip}).
\end{lemma}
\begin{proof}
This proof proceeds exactly as the proof of Lemma 4.4 for certifiable small-set expanders in~\cite{BBKSS}:
we apply Theorem~\ref{thm:structure-johnson} to each function $F_s$ and sum up over $s$. Doing so, we get the following inequality:
\begin{align}
\sum_s \ip{F_s, LF_s} \geq &(1 - \lambda_{r+1})(1- \gamma^{1/3}\exp(r))\sum_s \delta(F_s) - \sum_{j = 0}^r \frac{c_j \ell^j}{\gamma} \E_{a \in [n]^j}[\sum_{s}q_{a}(F_s)(\delta(F_s^{2}|_a) - \gamma)] \notag\\
&+ \frac{4}{3}\sum_s \E[(F_s^3 - F_s)\Pi_{\geq \lambda_r} F_{s}] \label{eq:sp-intermediate}
\end{align}

We now set $\gamma = \exp(-r)$ such that $(1-\gamma^{1/3}\exp(r)) = 1/2$. Further we bound each of the terms $\sum_s \ip{F_s, LF_s}, \sum_s \delta(F_s)$ and the Booleanity error: $\frac{4}{3}\sum_s \E[(F_s^3 - F_s)\Pi_{\geq \lambda_r} F_{s}]$.

First note that the fraction of low-valued vertices is small and in particular is at most $\frac{\viol(X)+\viol(X')}{1-\beta}$. Since the functions cover all of the high-valued vertices we get $\sum \delta(F_s) \gtrsim 1 - \frac{\viol(X) + \viol(X')}{1-\beta}$. Next, $\sum_s \ip{F_s, LF_s}$ counts the fraction of edges crossing the shift-partition. Every such edge must be violated by either $X$ or $X'$ or must be incident on a vertex with low value, therefore we get:  $\sum_s \ip{F_s, LF_s} \lesssim \viol(X)+\viol(X')+ \frac{2\viol(X)+2\viol(X')}{1-\beta}$.  We have put approximate inequalities here since there are some error terms generated because $F_s$'s are not exact indicator functions. 
The Booleanity term $B(F)$ in Theorem~\ref{thm:structure-johnson} is also small (after summing up) because $F_s$'s are approximate-indicators (note that it is $0$ for $0/1$ functions since $F^3 = F$)\footnote{Note that if the functions $F_s$ were defined using $\Ind[\val_u(X)\geq \beta]$ then we would obtain the lemma statement without the terms involving $\eta$ and with $\nu = 0$}. These statements have been made formal in the claims from Section 4.2 of~\cite{BBKSS}, specifically Claim 4.2, 4.7 and 4.9 therein. They are simple to prove given the properties of the approximate-indicator polynomial $p_{\beta,\nu}$ hence we omit them here. 

Using the above claims to bound each sum, plugging in the bounds in \eqref{eq:sp-intermediate} and rearranging we get the desired inequality.

\end{proof}

\subsubsection{Reducing Global Correlation}
The last ingredient we need for the proof of Lemma~\ref{lem:sp-j} is 
that given a pseudo-expectation of sufficiently high degree, one can 
construct different pseudo-expectations (which are conditionals 
of the initial pseudo-expectation) that have no global correlations.
More precisely:
\begin{lemma}\label{lem:ragh-tan-appln}
For all $\tau,p,\beta,\nu \in (0,1)$ and $|\Sigma|, D, D' \in \Z$ 
such that $D' \geq \max(D,\widetilde{\Omega}(1/\nu))$ the following holds.

Suppose there is a degree $D' + \widetilde{\Omega}(\frac{\log |\Sigma|}{p \nu \tau})$ pseudodistribution $\mu$ over UG assignments that satisfies $\cA_I$, and a polynomial $E(X,X')$ satisfying $\pE_{(X,X') \sim \mu \times \mu}[E(X,X')] \geq p$ and $\cA_I \vdash_D E(X,X') \leq 1$. Then for all subsets $S \subseteq [N]$, there exist subsets $A, B \subseteq S$ of size at most $O(\frac{\log|\Sigma|}{p\tau})$ and strings $y_A, y_B$ such that conditioning $\mu$ on the events $Y_A = y_A$ and $Y_B = y_B$ gives pseudodistributions $\mu_1$ and $\mu_2$ of degree at least $D'$ such that: 
\begin{enumerate}
\item $\pE_{X,X'\sim \mu_1\times\mu_2}[E(X,X')] \geq \frac{p}{2}$.
\item $\E_{\substack{u_1,v_1 \sim S: u_1 \neq v_1\\ u_2,v_2 \sim S: u_2 \neq v_2}}[I(Y_{u_1,v_1};Y_{u_2,v_2})] \leq \tau.$
\item $\E_{\substack{u_1,v_1 \sim S: u_1 \neq v_1\\ u_2,v_2 \sim S: u_2 \neq v_2}}[I(Y'_{u_1,v_1};Y'_{u_2,v_2})] \leq \tau$,
\end{enumerate}
where for $A \subseteq [N]$, $Y_{A} = (X_{i_1},\ldots, X_{i_{|A|}},p_{i_1},\ldots, p_{i_{|A|}})$ (same for $Y'$) and the mutual information is with respect to the collection of local distributions $\cL(\mu_1 \times \mu_2,p_{\beta,\nu})$.
\end{lemma}
\begin{proof}
Deferred to Section~\ref{sec:ragh}.
\end{proof}

\subsubsection{Proof of Lemma~\ref{lem:sp-j}}
We combine the lemmas stated in the previous section and set parameters to complete the proof of Lemma~\ref{lem:sp-j}.

\begin{lemma}[Restatement of Lemma~\ref{lem:sp-j}]\label{lem:sp-j-restated}
There exists a constant $\eps_0 \in (0,1)$, such that for all positive constants $\eps \leq \eps_0$, $\alpha, \tau \leq 1$, $\nu \leq \eps \exp(-r)$, all integers $\ell \geq \Omega(r), n \geq \ell$, where $r = \lfloor \frac{32\sqrt{\eps}}{\alpha}\rfloor$ the following holds: Let $I$ be an affine UG instance on $J_{n,\ell,\alpha}$ and $\mu$ be a pseudodistribution over assignments for $I$ with $\val_\mu(I) \geq 1-\eps$ and degree at least $\widetilde{\Omega}\left(\frac{|\Sigma|\ell^r\exp(r)}{\nu\tau\sqrt{\eps}}\right)$. There exists a restriction $a \subseteq [n]$ of size $j \leq r$, a degree $\widetilde{O}(1/\nu)$ polynomial $P_a(X,X')$ in a fixed set of $|\Sigma|$ polynomials, subsets $A, B \subseteq V(J|_a)$ of size at most $\tO(\frac{|\Sigma|\ell^j\exp(r)}{\tau\sqrt{\eps}})$ and strings $y_A, y_B$ such that conditioning $\mu$ on the events $Y_A = y_A$ and $Y_B = y_B$ gives degree $\widetilde{\Omega}(1/\nu)$ pseudodistributions $\mu_1$ and $\mu_2$ such that:
\begin{enumerate}
\item $\cA_I  ~~\vdash_{\widetilde{O}(1/\nu)} ~~ P_a(X,X') \in [0,1]$.
\item $\pE_{\mu_1 \times \mu_2}[\Phi^a_{\sqrt{\eps},\nu}(X,X')P_a(X,X')] \geq \exp(-r) \pE[P_a(X,X')]$.
\item $\pE_{\mu_1 \times \mu_2}[P_a(X,X')] \geq \Omega\left(\frac{\sqrt{\eps}}{|\Sigma|\ell^{j}\exp(r)}\right)$.
\item $\E_{\substack{u_1,v_1 \sim S: u_1 \neq v_1\\ u_2,v_2 \sim S: u_2 \neq v_2}}[I(Y_{u_1,v_1};Y_{u_2,v_2})] \leq \tau.$
\item $\E_{\substack{u_1,v_1 \sim S: u_1 \neq v_1\\ u_2,v_2 \sim S: u_2 \neq v_2}}[I(Y'_{u_1,v_1};Y'_{u_2,v_2})] \leq \tau,$
\end{enumerate}
where $S = V(J|_a)$, $Y_{u,v} = (X_u,X_v,p_u,p_v)$, $Y'_{u,v} = (X'_u,X'_v,p'_u,p'_v)$ and the mutual information is taken with respect to the collection of local distributions $\cL(\mu_1 \times \mu_2,p_{\sqrt{\eps},\nu})$. 
\end{lemma}
\begin{proof}
% We first prove that there exists a subcube $C$ such that the pseudoexpectation of the global shift-potential restricted to $C$ is large under an appropriate pseudodistribution $\mu_1 \times \mu_2$ when conditioned on a polynomial $P_C(X,X')$:
% \begin{equation}\label{eq:global-large-C}
% \pE_{\mu_1 \times \mu_2}[\Phi_{\beta,\nu}(X,X')|_C P_C(X,X')] \geq \exp(-r)\pE_{\mu_1 \times \mu_2}[P_C(X,X')],
% \end{equation}
% We need $\mu_1, \mu_2$ to satisfy certain low global correlation properties (points 4,5 of the lemma statement). Once we have this, it is easy to relate the global shift-potential on $C$ to the shift-potential on $C$ to get the above inequality albeit with $\Phi_{\beta,\nu}^C(X,X')$.

% Let us prove equation~\ref{eq:global-large-C}, while on the way obtaining $\mu_1,\mu_2$ with low global correlation properties. Recall that the functions $F_s$ can be thought of as (approximate) indicators of sets that form a partition of the graph that covers most vertices. Edges crossing this partition must be unsatisfied in either $X$ or $X'$., and applying Theorem~\ref{thm:structure-johnson} to each function $F_s$ and summing up over $s$ we get that:

We will apply Lemma~\ref{lem:structure-fs-j} with the following parameters: $r = \lfloor 32\sqrt{\eps}/\alpha \rfloor$, so that $1 - \lambda_{r+1} \geq 16\sqrt{\eps}$,  $\eta = 1/\sqrt{\eps}$ and $\beta = 201\sqrt{\eps}$.
% \footnote{Note that we have not optimized for these constants and chosen them large/small enough for ease of calculations.}
Since $\eps$ is sufficiently small and $\nu \leq \eps/2$ we get that $1/(1 - \beta - \nu) < 2$. By assumption $\pE_\mu[\viol(X)] = \pE_\mu[\viol(X')] \le \eps$, $\pE_\mu$ has degree $\widetilde{O}(1/\nu)$ and $\pE_\mu$ satisfies $\cA_I$, therefore taking the pseudoexpectation of equation~\ref{eq:sp-j-1} with respect to the pseudodistribution $\mu \times \mu$ we get:
\begin{align}
\sum_{j = 0}^r c_j \exp(r) \ell^j \E_{a \in {[n] \choose j}}\pE_{\mu \times \mu}[\sum_{s}q_{a}(F_s)(\delta(F_s|_a) - \exp(-r))] \geq \frac{\sqrt{\eps}}{3}.    
\end{align}

Since $c_j$'s are smaller than $\exp(r)$, an averaging argument gives us a size $j \leq r$ restriction $a \subseteq [n]$ and an $s \in \Sigma$ such that:
\begin{equation}\label{eq:sp-j-2}
\pE_{\mu \times \mu}[q_{a}(F_s)(\delta(F_s|_a) - \exp(-r))] \geq \Omega\left(\frac{\sqrt{\eps}}{|\Sigma|\ell^{j}\exp(r)}\right).
\end{equation}

Let $P_a(X,X')$ be the polynomial $q_a(F_s(X,X'))$ from above, abbreviated henceforth as $P_a$. First note that since $q_a$ is a degree $O(1)$-SoS polynomial we immediately get that $P_a$ is a degree $\widetilde{O}(1/\nu)$ SoS polynomial such that:
\[\cA_I ~~ \vdash_{\widetilde{O}(1/\nu)} F_s(u) \in [0,1] ~~\vdash_{O(1)} ~~ q_a(F_s(X,X')) = P_a(X,X') \in [0,1],\]
thus proving the first statement in the theorem.
We now reduce the global correlations using Lemma~\ref{lem:ragh-tan-appln}. 
Namely, we take $\mu$ with $S$ being the set of variables corresponding to the subcube $a$, the polynomial $E(X,X') = P_a(\delta(F_s|_a) - \exp(-r))$ (for which we know that $\cA \vdash_{\tO(1/\nu)} E(X,X') \leq 1$), $\beta = \sqrt{\eps}$ and $\tau,\nu$ as in the lemma statement. Thus, we get a 
pseudodistribution $\mu_1 \times \mu_2$ over UG assignments to the subcube $a$, such that: 
\begin{equation}\label{eq:sp-large-mu1-mu2}
\pE_{\mu_1 \times \mu_2}[P_a(\delta(F_s|_a) - \exp(-r))] \geq \Omega\left(\frac{\sqrt{\eps}}{|\Sigma|\ell^{j}\exp(r)}\right),
\end{equation}
 along with the conditions on the mutual information of $\mu_1,\mu_2$. 
%This step requires that $\deg(\mu) \geq \widetilde{\Omega}(\frac{|\Sigma|\log |\Sigma|\ell^j\exp(r)}{\nu\tau\sqrt{\eps}})$. 
This proves point (4),(5) from the lemma statement. 

\paragraph{Establishing point (2) of the lemma.}
Let us now derive the fact that the global shift-partition potential $\Phi(\mu_1 \times \mu_2)|_a$ is large conditioned on $P_a$.  
Rearranging~\eqref{eq:sp-large-mu1-mu2} we get that:
\begin{equation}\label{eq1}
    \pE_{\mu_1 \times \mu_2}[P_{a}\delta(F_s|_a)] \geq \exp(-r)\pE_{\mu_1 \times \mu_2}[P_a] + \Omega\left(\frac{\sqrt{\eps}}{|\Sigma|\ell^{j}\exp(r)}\right) \geq \exp(-r) \pE_{\mu_1 \times \mu_2}[P_a].
\end{equation}


Since $P_a$ is an SoS polynomial we can reweight $\pE_{\mu_1 \times \mu_2}[\cdot]$ by $P_a$ to get $\widetilde{\mathbb{E}'}$ and apply Cauchy-Schwarz to get:
\[\exp(-r) \leq \widetilde{\mathbb{E}'}[\delta(F_s|_a)] \leq \sqrt{\widetilde{\mathbb{E}'}[\delta(F_s|_a)^2]} = 
\sqrt{\frac{\pE_{\mu_1 \times \mu_2}[\delta(F_s|_a)^2 P_a]}{\pE_{\mu_1 \times \mu_2}[P_a]}},\]
rearranging which we get that: $\pE_{\mu_1 \times \mu_2}[\delta(F_s|_a)^2P_a] \geq \exp(-r)\pE_{\mu_1 \times \mu_2}[P_a]$. By definition, the global potential restricted to $a$ is equal to $\sum_{s \in \Sigma} \delta(F_s|_a)^2$. Therefore adding the terms $\pE_{\mu_1 \times \mu_2}[\delta(F_{s'}|_a)^2 P_a]$ for $s' \in \Sigma, s' \neq s$ (which are all non-negative) to the LHS of~\eqref{eq1} we get:
\begin{equation}\label{eq:global-large}
\pE_{\mu_1 \times \mu_2}[\Phi_{\beta,\nu}(X,X')|_a P_a] \geq \exp(-r)\pE_{\mu_1 \times \mu_2}[P_a].
\end{equation}

We will now relate the global potential to the shift-partition potential on $a$, where the only difference between the quantities is that in the former the value of a vertex is the fraction of edges satisfied in the whole graph, whereas in the latter it is the value calculated according to only the edges inside the subcube $a$ (in the terms $p_{\beta,\nu}(\val(u))$).

First note that that an $r$-restricted subcube has bounded expansion when $r$ is not too large: $\phi(J|_a) \leq O(\sqrt{\eps})$ (Claim~\ref{claim:subcube-expanse}). Using this we get that,
\[
\Phi^a_{\beta -200\sqrt{\eps},\nu}(X,X') \ge \Phi_{\beta,\nu}(X,X')|_a - 4\nu,
\]
which we prove formally in Claim~\ref{claim:potentials}.


% \begin{claim}\label{claim:subcube-expanse}
% If $r = \left\lfloor\frac{32\sqrt{\eps}}{\alpha}\right\rfloor < \frac{\ell}{4}$ and $s < r$, an $s$-restricted subcube of $J(n,\ell,\alpha)$ has expansion at most $200\sqrt{\eps}$.
% \end{claim}
% \begin{proof}
% A direct computation, deferred to Section~\ref{sec:aux_claims}.
% \end{proof}
% The proof of this claim is via a direct calculation, and we prove it later in this section. From this claim, we are able to prove that the local and global restricted potentials are related:

% \begin{claim}\label{claim:potentials}
% Suppose that $C$ is an $r$-restricted subcube of $J_{n,\ell,\alpha}$ with $r = \left\lfloor\frac{32\sqrt{\eps}}{\alpha}\right\rfloor$.
% Then if $\Phi^C$ is the shift-partition potential restricted to $C$, for any $\beta \ge 201\sqrt{\eps}$ and $\nu < \eps$,
% and furthermore this is certifiable in degree $\widetilde{O}(1/\nu)$ \sos.
% \end{claim}
% \begin{proof}
% A direct computation, deferred to Section~\ref{sec:aux_claims}.
% \end{proof}
% The proof of this claim is based on the fact that the fraction of neighbors of every vertex $v \in C$ which lie outside of $C$ cannot be too large when $r$ is bounded, and therefore if the value in $J$ at a vertex is $\beta$, the value restricted to $C$ is still $\beta - \phi(C)$. We prove it later in this section. This completes the proof of point (2) of the lemma, since plugging in the above claim into equation~\ref{eq:global-large} (with $\beta = 201\sqrt{\eps}$ set earlier in the proof) we get:
Overall, we get:
\[\pE_{\mu_1 \times \mu_2}[\Phi^a_{\sqrt{\eps},\nu}(X,X')P_a] \ge \pE_{\mu_1 \times \mu_2}[(\Phi_{201\sqrt{\eps},\nu}(X,X')|_a - 4\nu) P_a] \geq \exp(-r)\pE_{\mu_1 \times \mu_2}[P_a],\]
since $\nu \leq \exp(-r)/8$, which establishes point (2) of the 
lemma.

Point (3) follows from~\eqref{eq:sp-large-mu1-mu2}. Indeed, 
since $\delta(F_s|_a) \leq 1$ and $P_a(X,X') \geq 0$ (both facts are certifiable in degree $\widetilde{O}(1/\nu)$) we get that: $\cA_I \vdash_{\widetilde{O}(1/\nu)} P_{a}(X,X')(\delta(F_s|_a)-\exp(-r))] \leq (1-\exp(-r))P_a(X,X')$. Taking $\pE_{\mu_1 \times \mu_2}$ and plugging this into equation~\eqref{eq:sp-large-mu1-mu2} we get that:
\[\pE_{\mu_1 \times \mu_2}[P_a] \geq \Omega\left(\frac{\sqrt{\eps}}{|\Sigma|\ell^{j}\exp(r)}\right)
\qedhere\]
\end{proof}


\subsubsection{Auxiliary Claims}\label{sec:aux_claims}
We end this section by giving the proofs of Claims~\ref{claim:subcube-expanse} and \ref{claim:potentials}. These were also used in~\cite{BBKSS} but we include the proofs here for completeness.


\begin{claim}[Claim 6.10~\cite{BBKSS}]\label{claim:subcube-expanse}
If $r = \left\lfloor\frac{32\sqrt{\eps}}{\alpha}\right\rfloor < \frac{\ell}{4}$ and $s < r$, an $s$-restricted subcube of $J(n,\ell,\alpha\ell)$ has expansion at most $200\sqrt{\eps}$.
\end{claim}


\begin{proof}
Let $J|_Y$ be an $s$-restricted subcube.
We have that,
\[1 - \phi(J|_Y) = \frac{\binom{\ell - |Y|}{\alpha\ell}}{\binom{\ell}{\alpha\ell}} \geq \frac{\binom{\ell - r}{\alpha\ell}}{\binom{\ell}{\alpha\ell}} = \left(\frac{\ell - \alpha\ell}{\ell}\right)\left(\frac{\ell - \alpha\ell - 1}{\ell - 1}\right)\ldots \left(\frac{\ell - \alpha\ell - r + 1}{\ell - r + 1}\right).\]
Now since $r \leq \ell/4$ by assumption, each of the parenthesized terms is at least $\left(\frac{3\ell/4 - \alpha\ell}{3\ell/4}\right) = (1 - 4\alpha/3)$, so
\[1 - \phi(J|_Y) \geq \left(1-\frac{4\alpha}{3}\right)^r \geq 1 - \frac{4r\alpha}{3}.\]

Since $r = \left\lfloor\frac{32\sqrt{\eps}}{\alpha}\right\rfloor < \frac{75\sqrt{\eps}}{\alpha}$, we get that $\phi(J|_Y) < 200\sqrt{\eps}$
as desired.
\end{proof}


\begin{claim}[Claim 6.11~\cite{BBKSS}]\label{claim:potentials}
Suppose that $C$ is an $r$-restricted subcube of $J_{n,\ell,\alpha}$ with $r = \left\lfloor\frac{32\sqrt{\eps}}{\alpha}\right\rfloor$.
Then if $\Phi^C$ is the shift-partition potential restricted to $C$, for any $\beta \ge 201\sqrt{\eps}$ and $\nu < \eps$,
\[
\Phi^C_{\beta -200\sqrt{\eps},\nu}(X,X') \ge \Phi_{\beta,\nu}(X,X')|_C - 4\nu,
\]
and furthermore this is certifiable in degree $\tO(1/\nu)$ \sos.
\end{claim}


\begin{proof}
When $r = \left\lfloor\frac{32\sqrt{\eps}}{\alpha}\right\rfloor$, the expansion of $C$ is at most $1 - (1-4\alpha/3)^r \le 200\sqrt{\eps}$ by Claim~\ref{claim:subcube-expanse}.
Furthermore, from the definition of the Johnson graph this holds vertex-by-vertex; every $v \in C$ has at most a $200\sqrt{\eps}$-fraction of its neighbors outgoing.
Therefore,
\[
\Ind[\val_u^C(X) \ge \beta - 200\sqrt{\eps}] \ge \Ind[\val_u(X) \ge \beta],
\]
and furthermore since $\nu < \eps$,
\[
p_{\beta - 200 \sqrt{\eps},\nu}(\val_u^C(X)) + \nu \ge p_{\beta,\nu}(\val_u(X)) - \nu.
\]
The same holds for $p(\val_u(X'))$. Let $Z_{u,s} = \Ind[X_u - X'_u = s]$. Therefore, by definition,
\begin{align*}
    \Phi_{\beta - 200\sqrt{\eps},\nu}(X|_C,X'|_C)
    &= \sum_{s \in \Sigma} \E_{u \in C}\left[Z_{u,s}\cdot p_{\beta-200\sqrt{\eps},\nu}(\val_u^C(X))\cdot p_{\beta-200\sqrt{\eps},\nu}(\val_u^C(X'))\right]^2\\
    &\ge \sum_{s \in \Sigma} \E_{u \in C}\left[Z_{u,s}\cdot (p_{\beta,\nu}(\val_u(X)) - 2\nu)\cdot (p_{\beta,\nu}(\val_u(X')) - 2\nu)\right]^2\\
    &\ge \Phi_{\beta,\nu}(X,X')|_C - 4\nu,
\end{align*}
where each inequality is a sum-of-squares inequality of degree at most $2\deg(p)$.
\end{proof}








\subsection{Proof of Lemma~\ref{lem:ragh-tan-appln}: Reducing Global Correlation}\label{sec:ragh}
We will use the following lemma from~\cite{RT12}.

\begin{lemma}\label{lem:ragh-tan}
There exists $t \leq r$ such that: 
\[\E_{i_1,\ldots,i_t \sim [M]}\E_{i,j \sim [M]}[I(Y_i;Y_j \mid Y_{i_1}, \ldots,Y_{i_t})] \leq \frac{\log q}{r - 1},\]
where $q$ is the size of the domain of $Y_i$.
\end{lemma}

Note that the above lemma holds as long as there is a local collection of distributions over the variables $(Y_1,\ldots,Y_M)$ that are valid probability distributions over all collections of $t+2$ variables and are consistent with each other. We will apply the above lemma to reduce the global correlation between $(Y_{u_1,v_1};Y_{u_2,v_2})$ (similarly for $Y'$), for the local collection of distributions $\cL$ (Definition~\ref{defn:local-dist}).


\begin{lemma}[Lemma~\ref{lem:ragh-tan-appln} restated]
For all $\tau,p,\beta,\nu \in (0,1)$ and $|\Sigma|, D, D' \in \Z$, $D' \geq \max(D,\widetilde{\Omega}(1/\nu))$ the following holds: Suppose there is a degree $D' + \widetilde{\Omega}(\frac{\log |\Sigma|}{p \nu \tau})$ pseudodistribution $\mu$ over UG assignments that satisfies $\cA_I$, and a polynomial $E(X,X')$ satisfying $\pE_{(X,X') \sim \mu \times \mu}[E(X,X')] \geq p$ and $\cA_I \vdash_D E(X,X') \leq 1$. Then for all subsets $S \subseteq [N]$, there exist subsets $A, B \subseteq S$ of size at most $O(\frac{\log|\Sigma|}{p\tau})$ and strings $y_A, y_B$ such that conditioning $\mu$ on the events $Y_A = y_A$ and $Y_B = y_B$ gives pseudodistributions $\mu_1$ and $\mu_2$ of degree at least $D'$ such that: 
\begin{enumerate}
\item $\pE_{X,X'\sim \mu_1\times\mu_2}[E(X,X')] \geq \frac{p}{2}$.
\item $\E_{\substack{u_1,v_1 \sim S: u_1 \neq v_1\\ u_2,v_2 \sim S: u_2 \neq v_2}}[I(Y_{u_1,v_1};Y_{u_2,v_2})] \leq \tau.$
\item $\E_{\substack{u_1,v_1 \sim S: u_1 \neq v_1\\ u_2,v_2 \sim S: u_2 \neq v_2}}[I(Y'_{u_1,v_1};Y'_{u_2,v_2})] \leq \tau$,
\end{enumerate}
where for $A \subseteq [N]$, $Y_{A} = (X_{i_1},\ldots, X_{i_{|A|}},p_{i_1},\ldots, p_{i_{|A|}})$ (same for $Y'$) and the mutual information is with respect to the collection of local distributions $\cL(\mu_1 \times \mu_2,p_{\beta,\nu})$.
\end{lemma}


\begin{proof}
Without loss of generality suppose that $S \subseteq [N]$ is $[m]$. Let ${[m] \choose r}$ denote the set of $r$-sized subsets of $[m]$ and $m_r$ denote ${m \choose r}$. Let $Y = (Y_{a_1},\ldots,Y_{a_{m_2}})$ for $a_t = (u,v) \in {[m] \choose 2}$. For brevity we will use $[m_2]$ for ${[m] \choose 2}$. Given the pseudodistribution $\mu \times \mu$ over assignments $(X,X')$ we will consider the 
local collection of distributions $\cL(\mu \times \mu,p_{\beta,\nu})$ over the variables of $X,X',p,p'$ and let this also denote the induced collection of distributions over $(Y,Y')$ (Definition~\ref{defn:local-dist}). Recall the notation for random variable $Y_A$, for any $A \subseteq [m]$: $Y_A = (X_{i_1}, \ldots, X_{i_{|A|}}, p_{i_1}, \ldots, p_{i_{|A|}})$. We will use $y_S$ to denote instantiations of $Y_S$.

Let $\cL$ denote $\cL(\mu \times \mu,p_{\beta,\nu})$.
Applying Lemma~\ref{lem:ragh-tan} to the collection of  distributions $\cL$ over $Y$ we get that for all $\tau' > 0$, conditioning on $t = O(\frac{\log |\Sigma|}{\tau'})$ variables gives:
\begin{equation}\label{eq:rt-1}
\E_{i_1,\ldots,i_t \sim [m_2]}[\E_{a,b \sim [m_2]} I(Y_a ; Y_b \mid Y_{i_1}, \ldots,Y_{i_t})] \leq \tau'.
\end{equation}

Let $\cL|_R$ denote the distribution that $\cL$ induces on the random variable $R$. We know that the distribution $\cL|_{(Y_{i_1}, \ldots, Y_{i_t})}$ can be reduced to the distribution $\cL|_{Y_{i_1 \cup \ldots \cup i_t}}$ by discarding repeating indices in $i_1,\ldots,i_t$. This is because $\cL$ is consistent on the value it assigns to a vertex $u$ when it occurs in $Y_{u,v}$ or $Y_{u,v'}$. We will drop the subscript $R$ from $\cL|_R$ when the random variable is clear from context. Expanding the definition of conditional mutual information we therefore get:
\begin{align*}
I(Y_a;Y_b \mid Y_{i_1}, \ldots, Y_{i_t}) &= \E_{(y_{i_1}, \ldots,y_{i_t}) \sim \cL|_{(Y_{i_1}, \ldots, Y_{i_t})}}[I(Y_a ; Y_b \mid Y_{i_1} = y_{i_1}, \ldots,Y_{i_t} = y_{i_t})] \\
&= \E_{y_{i_1 \cup \ldots \cup i_t} \sim \cL|_{Y_{i_1 \cup \ldots \cup i_t}}}[I(Y_a ; Y_b \mid Y_{i_1 \cup \ldots \cup i_t} = y_{i_1 \cup \ldots \cup i_t}].
\end{align*}
Plugging the above into equation~\eqref{eq:rt-1} and applying Markov's inequality we get that for all $\alpha \in (0,1)$:
\begin{equation}\label{eq:avg1}
\Pr_{\substack{i_1,\ldots,i_t \sim [m_2] \\ y_{i_1 \cup \ldots \cup i_t} \sim \cL}}\left[\E_{a,b \sim [m_2]}[I(Y_a ; Y_b \mid Y_{i_1 \cup \ldots \cup i_t} = y_{i_1 \cup \ldots \cup i_t})] \geq \frac{\tau'}{\alpha} \right] \leq \alpha.
\end{equation}
Repeating the above analysis for $Y'$ we get:
\begin{equation}\label{eq:avg2}
\Pr_{\substack{j_1,\ldots,j_t \sim [m_2] \\ y'_{j_1 \cup \ldots \cup j_t} \sim \cL}}\left[\E_{a,b \sim [m_2]}[I(Y_a ; Y_b \mid Y'_{j_1 \cup \ldots \cup j_t} = y_{j_1 \cup \ldots \cup j_t})] \geq \frac{\tau'}{\alpha} \right] \leq \alpha.
\end{equation}

Now recall that (Definition~\ref{defn:local-dist}) the probability of the event $(X_u,p_u) =  (\sigma,b)$ under $\cL$, for some $\sigma \in \Sigma$ and $b \in \{0,1\}$, corresponds to the pseudoexpectation under $\mu$ of an appropriate polynomial we will denote by $Q_{u,(\sigma,b)}$:
\begin{align}\label{eq:poly-cond}
\pPr_{\cL}[(X_u, p_u) = (\sigma, b)] &= \pE[X_{u,\sigma}(1-b + (-1)^{1-b}p(\val_u(X)))] \nonumber \\
&:= \pE_\mu[Q_{u,(\sigma,b)}].
\end{align}
We can extend the above definition by letting $Q_{S, \cup_{u \in S} (\sigma_u,b_u)}$ denote the polynomial $\prod_{u \in S} Q_{u,(\sigma_u,b_u)}$. For $y_S = (\cup_{u \in S} \sigma_u, \cup_{u \in S} b_u)$ we get that:
\begin{align}\label{eq:poly-cond-prod}
\pPr_{\cL}[(\cup_{u \in S} X_u, \cup_{u \in S} p_u) = Y_S = y_S] = \pE_\mu[Q_{S,\cup_{u \in S}(\sigma_u,b_u)}].
\end{align}
One can define analogous notation for $Y'$ (e.g. the polynomials $Q'_{u,(\sigma,b)}$, etc) and derive the statements above. We know that $\cA_I \vdash_{t~\tO(1/\nu)} Q_{S, \cup_{u \in S} (\sigma_u,b_u)} \in [0,1]$. For ease of notation let us define the following expression:
\begin{align*}
\pE_{\mu \times \mu}[E(X,X') \mid Y_S = y_S, Y'_T = y'_T] &:= \pE_{\mu \times \mu}[E(X,X') \mid  Q_{S, y_S} \cdot  Q'_{T, y'_T}] \\
&= \frac{\pE_{\mu \times \mu}[E(X,X')Q_{S, y_S} Q'_{T, y'_T}]}{\pE_{\mu \times \mu}[Q_{S, y_S} Q'_{T, y'_T}]} \\
&= \frac{\pE_{\mu \times \mu}[E(X,X')Q_{S, y_S} Q'_{T, y'_T}]}{\pPr_{\cL}[Y_{S} = y_{S}, Y'_{T} = y'_{T}]},
\end{align*}
where the last equality follows from equation~\eqref{eq:poly-cond-prod}. Analogous to the definition of conditional expectation we can check that:
\[\pE_{\mu \times \mu}[E(X,X')] = \E_{\substack{i_1,\ldots,i_t \sim [m] \\ j_1,\ldots,j_t \sim [m] \\ y_{i_1 \cup \ldots \cup i_t}, y'_{j_1 \cup \ldots \cup j_t} \sim \cL}}[\pE_{X,X'}[E(X,X') \mid Y_{i_1 \cup \ldots \cup i_t} = y_{i_1 \cup \ldots \cup i_t}, Y'_{j_1 \cup \ldots \cup j_t} = y'_{j_1 \cup \ldots \cup j_t}]].\]
%where $\Pr_{X \sim D}[E(X) \mid Y]$ is defined as $\E_{y \sim D}[\Pr_{X \sim D}[E(X) \mid Y = y]]$.

Since $\cA_I \vdash_D E(X,X') \leq 1$ and $\deg(\mu) \geq D + \widetilde{\Omega}(t/\nu)$ we get that $E(X,X') \leq 1$ even after conditioning on a non-negative event $Q$ of degree $\widetilde{O}(t/\nu)$: $\pE[E(X,X') \mid Q] \leq 1$. An averaging argument implies that:
\begin{equation}\label{eq:avg3}
\Pr_{\substack{i_1,\ldots, j_1,\ldots \sim [m_2] \\ y_{i_1 \cup \ldots \cup i_t},y'_{j_1 \cup \ldots \cup j_t} \sim \cL}}\left[\pE_{\mu \times \mu}[E(X,X') \mid Y_{i_1 \cup \ldots \cup i_t} = y_{i_1 \cup \ldots \cup i_t}, Y'_{j_1 \cup \ldots \cup j_t} = y'_{j_1 \cup \ldots \cup j_t}] \geq \frac{p}{2}\right] \geq \frac{p}{2}.
\end{equation}

Choosing $\alpha = p/8$ and $\tau' < p\tau/8$ we can take a union bound over the events in equations~\eqref{eq:avg1},~\eqref{eq:avg2},~\eqref{eq:avg3} to get that there exist sets $i_1,\ldots,i_t, j_1,\ldots,j_t \in [m_2]$ and $y_{i_1 \cup \ldots \cup i_t}, y'_{j_1 \cup \ldots \cup j_t} \sim \cL$ such that,   
\[\E_{a,b \sim [m_2]}[I(Y_a ; Y_b \mid Y_{i_1 \cup \ldots \cup i_t} = y_{i_1 \cup \ldots \cup i_t}] \leq \frac{\tau'}{\alpha} < \tau.\]
\[\E_{a,b \sim [m_2]}[I(Y'_a ; Y'_b \mid Y'_{j_1 \cup \ldots \cup j_t} = y'_{j_1 \cup \ldots \cup j_t})] \leq \frac{\tau'}{\alpha} < \tau.
\]
\[\pE_{\mu \times \mu}[E(X,X') \mid Y_{i_1 \cup \ldots \cup i_t} = y_{i_1 \cup \ldots \cup i_t}, Y'_{j_1 \cup \ldots \cup j_t} = y'_{j_1 \cup \ldots \cup j_t}] \geq \frac{p}{2}.\]

Let $\mu_1$ be the pseudodistribution on $X$ that we get by conditioning $\mu$ on $Q_{i_1 \cup \ldots \cup i_t, y_{i_1 \cup \ldots \cup i_t}}$ and let $\mu_2$ be the pseudodistribution on $X'$ that we get by conditioning $\mu$ on $Q_{j_1 \cup \ldots \cup j_t,y'_{j_1 \cup \ldots \cup j_t}}$. It is easy to check that:
\[\cL(\mu \times \mu, p_{\beta,\nu})| (Y_{i_1 \cup \ldots \cup i_t} = y_{i_1 \cup \ldots \cup i_t}, Y'_{j_1 \cup \ldots \cup j_t} = y'_{j_1 \cup \ldots \cup j_t}) = \cL(\mu_1 \times \mu_2, p_{\beta,\nu}),\]
thus giving us the three properties we need in the lemma. It remains to check the degree bounds on $\mu$ that we require. We have conditioned on $t = O(\log |\Sigma|/p\tau)$ variables from $(Y,Y')$, each having degree $\tO(1/\nu)$ in $X,X'$, therefore it suffices to have degree of $\mu$ to be $\widetilde{\Omega}\left(\frac{\log |\Sigma|}{p\tau\nu}\right)$. After conditioning we get that $\deg(\mu_1 \times \mu_2) \geq D'$ if $\deg(\mu) \geq D'+\widetilde{\Omega}\left(\frac{\log |\Sigma|}{p\tau\nu}\right)$ thus completing the proof.
\end{proof}

% and $Y'$ separately, 
% \[\E_{\substack{j_1,\ldots,j_t \sim [m_2] \\ (y'_{j_1},\ldots,y'_{j_t}) \sim \cL}}[\E_{a,b \sim [m_2]} I(Y'_a ; Y'_b \mid Y'_{j_1} = y'_{j_1}, \ldots,Y'_{j_t} = y'_{j_t})] = \E_{j_1,\ldots,j_t \sim [m_2]}[\E_{a,b \sim [m_2]} I(Y'_a ; Y'_b \mid Y'_{j_1}, \ldots,Y'_{j_t})] \leq \tau'.\]

% \begin{equation}\label{eq:avg2}
% \Pr_{\substack{j_1,\ldots,j_t \sim [m_2] \\ (y'_{j_1},\ldots,y'_{j_t}) \sim \cL}}\left[\E_{a,b \sim [m_2]}[I(Y'_a ; Y'_b \mid Y_{j_1} = y'_{j_1}, \ldots,Y_{j_t}=y'_{j_t})] \geq \frac{\tau'}{\alpha} \right] \leq \alpha.
% \end{equation}

% Since we finally want to be able to relate $\cL$ to $\mu$ it will be useful to consider the probabilities over $i_t$'s that do not intersect. Modifying equations~\ref{eq:avg1},~\ref{eq:avg2} we get:
% \begin{equation}\label{eq:avg1}
% \Pr_{\substack{(i_1,\ldots,i_t) \sim [{m \choose 2t}] \\ (y_{i_1},\ldots,y_{i_t}) \sim \cL}}\left[\E_{a,b \sim [m_2]}[I(Y_a ; Y_b \mid Y_{i_1} = y_{i_1}, \ldots,Y_{i_t}=y_{i_t})] \geq \frac{\tau'}{\alpha} \right] \leq \alpha + O\left(\frac{t^2}{m^2}\right).
% \end{equation}
% \begin{equation}\label{eq:avg2}
% \Pr_{\substack{(j_1,\ldots,j_t) \sim {[m] \choose 2t} \\ (y'_{j_1},\ldots,y'_{j_t}) \sim \cL}}\left[\E_{a,b \sim [M]}[I(Y'_a ; Y'_b \mid Y_{j_1} = y'_{j_1}, \ldots,Y_{j_t}=y'_{j_t})] \geq \frac{\tau'}{\alpha} \right] \leq \alpha + O\left(\frac{t^2}{m^2}\right),
% \end{equation}
% where $(i_1,\ldots,i_t) \sim [{[m] \choose 2t}]$ refers to the distribution where we first draw a random $2t$-sized subset of $m$, and then randomly divide it into $2$-sized subsets $i_1,\ldots,i_t$.

% When the sets $i_1 = (u_1,v_1), \ldots, i_t = (u_t,v_t)$ from equations~\ref{eq:avg1} (similarly for the $j_t$'s) are disjoint then we have that\footnote{Note that the equation is not necessarily true when the sets are not disjoint. For e.g. if $i_1 = i_2$, then $\pPr_{\cL}[Y_{i_1} = Y_{i_2} = y] = \pE_\mu[Q_y]$, which may not be equal to $\pE_{\mu}[Q_y^2]$. This arises since the variables $p_u$ are not perfect indicators in the sense that $\pE_\mu[p_u]$ might not be equal to $\pE_{\mu}[p_u^2]$.}:
% \begin{align}\label{eq:y-event-conj}
% \pPr_{\cL}[Y_{i_1} = y_1, \ldots, Y_{i_t} = y_t] &= \pE_\mu[Q_{i_1,y_1} \cdot \ldots \cdot Q_{i_t,y_t}].
% \end{align}


\subsection{Proof of Lemma~\ref{lem:correlation}: 
Conditioning Does Not Introduce Correlations}

In this section we prove that given a product distribution over $(X,X')$ that satisfies low global correlation with respect to the variables $(Y,Y')$, where $Y_{u,v} = (X_u,X_v,p_u,p_v)$, and an event $E(X,X')$ that holds with large probability, conditioning on $E$ cannot correlate too many pairs $(Y_i,Y'_i)$. This is a general lemma that holds if the variables $(Y,Y')$ are local functions of the underlying variables $(X,X')$ and satisfy low global correlation, but here we state it for our specific application only. 


\begin{lemma}[Restatement of Lemma~\ref{lem:correlation}]
For all $\tau,p,\beta,\nu,\delta \in (0,1)$, integers $D,N$, $m \leq N$ the following holds: Let $\mu_1 \times \mu_2$ be a degree $D+\widetilde{\Omega}(1/\nu)$ pseudodistribution over $(X,X')$ satisfying $\cA_I(X,X')$ and $E(X,X')$ be a polynomial such that, $\cA_I \vdash_D E(X,X') \in [0,1]$ and $\pE_{\mu_1 \times \mu_2}[E(X,X')] \geq p$. Suppose for $S \subseteq [N]$ we have that, 
\begin{enumerate}
\item $\E_{\substack{u_1,v_1 \sim S: u_1 \neq v_1\\ u_2,v_2 \sim S: u_2 \neq v_2}}[I(Y_{u_1,v_1};Y_{u_2,v_2})] \leq \tau.$
\item $\E_{\substack{u_1,v_1 \sim S: u_1 \neq v_1\\ u_2,v_2 \sim S: u_2 \neq v_2}}[I(Y'_{u_1,v_1};Y'_{u_2,v_2})] \leq \tau$,
\end{enumerate}
where $Y_{u,v} = (X_u,X_v,p_u,p_v), Y'_{u,v} = (X'_u,X'_v,p'_u,p'_v)$ and the mutual information is with respect to the collection of local distributions $\cL(\mu_1 \times \mu_2,p_{\beta,\nu})$ (Definition~\ref{defn:local-dist}). Then we have that, 
\[\Pr_{u,v \sim S: u \neq v}[ TV((Y_{u,v},Y'_{u,v})|E , (Y_{u,v},Y'_{u,v})) \geq \delta] \leq O\left(\frac{\sqrt{\tau}+ 1/|S|}{p\delta^2}\right),\]
where the distribution $(Y_{u,v},Y'_{u,v})|E$ refers to the joint distribution on these variables defined by the collection of local distributions $\cL(\mu_1 \times \mu_2|E, p_{\beta,\nu})$ and similarly $(Y_{u,v},Y'_{u,v})$ refers to the distribution defined by $\cL(\mu_1 \times \mu_2,p_{\beta,\nu})$.
\end{lemma}

\begin{proof}
As in the proof of Lemma~\ref{lem:ragh-tan-appln}, without loss of generality suppose $S = [m]$, let $[m_2]$ denote $2$-sized subsets of $m$ and $Y = (Y_{a_1},\ldots,Y_{a_M})$ for $a_t = (u,v), (u,v) \in [m_2]$. Given the pseudodistribution $\cD$ over assignments $(X,X')$ we will consider the local collection of distributions $\cL(\cD, p_{\beta,\nu})$ over the variables $X,X',p,p'$ and therefore the induced collection of distributions over $(Y,Y') = (Y_{a_1},\ldots,Y_{a_M},Y'_{a_1},\ldots,Y'_{a_M})$ (Definition~\ref{defn:local-dist}).

Throughout the proof we will consider random variables $(Y,Y')$ drawn from the collection of local distributions $\cL(\mu_1 \times \mu_2,p_{\beta,\nu})$ and from the conditioned collection of distributions $\cL(\mu_1 \times \mu_2 | E,p_{\beta,\nu})$. For brevity of notation we will use $(Y_a,Y'_a)|E$ to be the joint distribution induced by $\cL(\mu_1 \times \mu_2 | E, p_{\beta,\nu})$ on $(Y_a,Y'_a)$ and use $(Y_a,Y'_a)$ to be the distribution induced by $\cL(\mu_1 \times \mu_2,p_{\beta,\nu})$. We will abbreviate $\cL(\mu_1 \times \mu_2,p_{\beta,\nu})$ to $\cL$. Let $U$ be the set of variables $a \in [m_2]$ for which $TV((Y_a,Y'_a)|E , (Y_a,Y'_a)) \geq \delta$ and let the fractional-size of $U$ be $\gamma$. If $\Omega$ denotes the domain of $Y_u, Y'_u$ then for every $a \in U$ there exists a set $S_a \subseteq \Omega \times \Omega$ such that: 
\begin{equation}\label{eq:e-tv}
\pPr_{\cL}[(Y_a,Y'_a) \in S_a \mid E] - \pPr_{\cL}[(Y_a,Y'_a) \in S_a] \geq \delta.
\end{equation}
Let $e_a$ denote  $\pPr_{\cL}[(Y_a,Y'_a) \in S_a]$. Define the random variables $Z_a = \one((Y_a,Y'_a) \in S_a) - e_a$. Define:
\[Z = \E_{a \sim U}[Z_a] = \E_{a \in U}[\one((Y_a,Y'_a) \in S_a) - e_a].\]

Let $\pE_{\cL}$ denote the natural expectation operator corresponding to the local distributions $\cL$. 
One can check that $\pE_{\cL}[Z] = \pE_{\cL}[Z_a] = 0$, and we now calculate its variance. For two events $A,B$ on the variables $Y,Y'$ let $\pCov(A,B)$ denote $\pE_{\cL}[AB] - \pE_{\cL}[A]\pE_{\cL}[B]$. Firstly for all $a,b \in [m_2]$  using Pinsker's inequality and data processing inequality we have that,
\begin{align*}
\pCov(\one((Y_a,Y'_a) \in S_a),\one((Y_b,Y'_b) \in S_b)) 
&\leq TV((Y_a,Y'_a,Y_b,Y'_b), (Y_a,Y'_a)\times (Y_b,Y'_b)) \\ &\leq O(\sqrt{I(Y_a,Y'_a;Y_b,Y'_b)}) \\
&\leq O(\sqrt{I(Y_a;Y_b) + I(Y'_a;Y'_b)})
\end{align*}

The proof will proceed by proving upper and lower bounds on $\pE_{\cL}[Z^2]$, where the upper bound uses low global correlation properties of $Y$ and $Y'$ and the lower bound uses the large deviation we have by equation~\ref{eq:e-tv}. 
\paragraph{Upper bound for $\pE_{\cL}[Z^2]$:}
We have the following upper bound:
\begin{align*}
\pE_{\cL}[Z^2] &= \E_{a,b \sim U}[\pE_{\cL}[(\one((Y_a,Y'_a) \in S_a) - e_a)(\one((Y_b,Y'_b) \in S_b) - e_b)]] \\
&= \E_{a,b \sim U}[\pE_{\cL}[\one((Y_a,Y'_a) \in S_a) - e_a]\pE_{\cL}[\one((Y_b,Y'_b) \in S_b) - e_b] + \pCov(\one((Y_a,Y'_a) \in S_a),\one((Y_b,Y'_b) \in S_b))] \\
&\leq \E_{a,b \sim U}[O(\sqrt{I(Y_a;Y_b) + I(Y'_a;Y'_b)})] \\
&\leq O(\sqrt{\E_{a,b \sim U}[I(Y_a;Y_b) + I(Y'_a;Y'_b)}])\\
&\leq O\left(\frac{\sqrt{\tau}}{\gamma}\right),
\end{align*}
where the last inequality follows because $\E_{a,b \sim [m_2]}[I(Y_a;Y_b)] \geq \gamma^2 \E_{a,b \sim U}[I(Y_a;Y_b)]$, and by assumption $\E_{a,b \sim [M]}[I(Y_a;Y_b)] \leq \tau$ (similarly for $Y'$).

\paragraph{Lower bound for $\pE_{\cL}[Z^2]$:}
Let $Q_a$ be the polynomial corresponding to the event $\Ind[(Y_a,Y'_a) \in S_a]$, i.e. $\pPr_{\cL}[(Y_a,Y'_a) \in S_a] = \pE_{\mu_1 \times \mu_2}[Q_a]$. Let $P(X,X')$ abbreviated as $P$ denote the polynomial $\E_a[Q_a - e_a]$. Let $Q_{a \cup b}$ be the polynomial such that $\pPr_{\cL}[(Y_a,Y'_a) \in S_a \wedge (Y_b,Y'_b) \in S_b] = \pE_{\mu_1 \times \mu_2}[Q_{a \cup b}]$. If $a \cap b = \phi$ we have that $Q_{a \cup b} = Q_a \cdot Q_b$, whereas this may not hold if they intersect in one or two variables. But we have that $\Pr_{a, b \sim U}[a \cap b \neq \phi] \leq O(1/|U|)$.
Using these facts we first show that, $\pE_{\cL}[Z^2]+O(1/|U|) \geq \pE_{\mu_1 \times \mu_2}[P^2]$ via the following two equations:
\begin{align}
\pE_{\cL}[Z^2] &= \E_{a,b \sim U}[\pE_{\cL}[(\one((Y_a,Y'_a) \in S_a) - e_a)(\one((Y_b,Y'_b) \in S_b) - e_b)]] \notag \\
&=\Pr_{a,b \sim U}[a \cap b = \phi]\E_{a,b \sim U: a\cap b=\phi}[\pE_{\mu_1 \times \mu_2}[(Q_{a} - e_a)(Q_b - e_b)]] \notag\\ &+ \Pr_{a,b \sim U}[a \cap b \neq \phi]\E_{a,b \sim U: a\cap b\neq \phi}[\pE_{\cL}[(\one((Y_a,Y'_a) \in S_a) - e_a)(\one((Y_b,Y'_b) \in S_b) - e_b)]] \notag \\
&\geq \Pr_{a,b \sim U}[a \cap b = \phi]\E_{a,b \sim U: a\cap b=\phi}[\pE_{\mu_1 \times \mu_2}[(Q_{a} - e_a)(Q_b - e_b)]] - O\left(\frac{1}{|U|}\right).\label{eq:z2-1}
\end{align}

On the other hand we have that:
\begin{align}
\pE_{\mu_1 \times \mu_2}[P^2] &= \E_{a,b \sim U}[\pE_{\mu_1 \times \mu_2}[(Q_a - e_a)(Q_b - e_b)]] \notag \\
&= \Pr_{a,b \sim U}[a \cap b = \phi]\E_{a,b \sim U: a\cap b=\phi}[\pE_{\mu_1 \times \mu_2}[(Q_{a} - e_a)(Q_b - e_b)]] \notag \\
&+\Pr_{a,b \sim U}[a \cap b \neq \phi]\E_{a,b \sim U: a\cap b=\phi}[\pE_{\mu_1 \times \mu_2}[(Q_{a} - e_a)(Q_b - e_b)]] \notag \\
&\leq \Pr_{a,b \sim U}[a \cap b = \phi]\E_{a,b \sim U: a\cap b\neq\phi}[\pE_{\mu_1 \times \mu_2}[(Q_{a} - e_a)(Q_b - e_b)]] + O\left(\frac{1}{|U|}\right)\label{eq:z2-2}
\end{align}
Combining \eqref{eq:z2-1} and \eqref{eq:z2-2} we get that:
\[\pE_{\cL}[Z^2]+O\left(\frac{1}{|U|}\right) \geq \pE_{\mu_1 \times \mu_2}[P^2].\]

Recall that $|U| = \gamma |S|$. Since $\cA_I \vdash_{D} E(X,X') \in [0,1]$ we get that:
\[\pE_{\cL}[Z^2]+O\left(\frac{1}{\gamma|S|}\right)  \geq \pE_{\mu_1 \times \mu_2}[P^2] \geq \pE_{\mu_1 \times \mu_2}[E]\pE_{\mu_1 \times \mu_2}[P^2 | E] \geq p \pE_{\mu_1 \times \mu_2}[P | E]^2 = p\pE_{\cL}[Z | E]^2,\]
where the last inequality is by Cauchy-Schwarz on the pseudodistribution $\mu_1 \times \mu_2 | E$ and the last equality follows by the definition of $P$ and $Z$. Using equation~\eqref{eq:e-tv} we know that for all $a \in U$:
\[\pE_{\cL}[Z_a \mid E] = \pPr_{\cL}[(Y_a,Y'_a) \in S_a \mid E] - e_a  \geq \delta,\]
which implies that $\pE_{\cL}[Z \mid E] \geq \delta$.

Combining the upper and lower bounds on $\pE_{\cL}[Z^2]$ we get that $\gamma \leq O\left(\frac{\sqrt{\tau}+ 1/|S|}{p\delta^2}\right)$,
completing the proof of the lemma.


%\mitali{SoS-ize}

%Applying Chebyshev's inequality we get that, for all $t > 0$,
%\[\Pr_{\mu_1 \times \mu_2}[Z \geq \frac{t\tau^{1/4}}{\gamma^{1/2}}] \leq \frac{\gamma\E_{Y,Y'}[Z^2]}{t^2 \sqrt{\tau}} \leq  O\left(\frac{1}{t^2}\right).\]


% Since $Z$ is always smaller than $1$, we can apply an averaging argument to get that,
% \[\Pr[Z \geq \frac{t\tau^{1/4}}{\gamma^{1/2}} \mid E] \geq \frac{\delta - \frac{t\tau^{1/4}}{\gamma^{1/2}}}{1 - \frac{t\tau^{1/4}}{\gamma^{1/2}}}.\]
% Since $E$ occurs with probability at least $p$ we get that, 
% \[\Pr[Z \geq \frac{t\tau^{1/4}}{\gamma^{1/2}}] \geq \frac{p(\delta - \frac{t\tau^{1/4}}{\gamma^{1/2}})}{1 - \frac{t\tau^{1/4}}{\gamma^{1/2}}}.\]
% Setting $t = \frac{\delta\gamma^{1/2}}{2\tau^{1/4}}$ and combining the upper and lower bounds we get that,
% \[\frac{p(\delta - \frac{t\tau^{1/4}}{\gamma^{1/2}})}{1 - \frac{t\tau^{1/4}}{\gamma^{1/2}}} = \frac{p\delta}{2 - \delta} \leq O\left(\frac{\sqrt{\tau}}{\delta^2\gamma}\right),\]
% rearranging which gives us the needed lower bound:
% \[\gamma \leq O\left(\frac{\sqrt{\tau}}{p\delta^3}\right).\]

%{\color{red} do we need to say that the RHS probability is $< 1$, in which case we're already done.}

\end{proof}




